% Write-up plan for the whole GT2027 paper: figure, table and equation lists, then one section per paper
% section (subsection headings, content, figures, tables, equations and their purpose).
% Build: bash build_plan.sh (pdflatex twice; --figs also redraws the figures).
% Paper figures are read from ASME_2027/figures/, each drawn by figures/code/<same name>.py.
\documentclass[10pt]{article}
\usepackage[utf8]{inputenc}
\usepackage[T1]{fontenc}
\usepackage{lmodern}
\usepackage{textcomp}
\usepackage[a4paper,margin=18mm]{geometry}
\usepackage{amsmath,amssymb,bm}
\usepackage{booktabs,float,graphicx,array,enumitem,tabularx}
\usepackage[dvipsnames]{xcolor}
\usepackage[colorlinks=true,linkcolor=MidnightBlue,urlcolor=MidnightBlue]{hyperref}
\setlength{\parskip}{4pt}
\setlength{\parindent}{0pt}
\setlist{itemsep=1pt,topsep=2pt}
\setlength{\emergencystretch}{2em}
\newcommand{\um}{\textmu m}
\graphicspath{{../figures/}}
\newcommand{\deta}{\Delta_{\mathrm{mesh}}/\eta_{K,193}}
% candidate paper figure: file in ../figures, width fraction, caption; a box if the file is not drawn yet
\newcommand{\planfig}[3]{\begin{figure}[H]\centering
  \IfFileExists{../figures/#1}{\includegraphics[width=#2\textwidth]{#1}}{\fbox{\parbox{0.6\textwidth}{\centering
  \texttt{\detokenize{#1}}: not drawn yet}}}\caption{#3}\end{figure}}

\title{GT2027 paper: write-up plan}
\author{}
\date{}

\begin{document}
\maketitle
\vspace{-14mm}

The plan for the whole paper, section by section: subsection headings, content, and the figures, tables and
equations with their purpose. Sections: 1 Introduction, 2 Research Aims, 3 Methods, 4 Results (4.1 Coolant Flow,
4.2 Breakdown of Model Differences, 4.3 A Posteriori Tests) and 5 Conclusions. Figures are in
\texttt{ASME\_2027/figures/}, each drawn by the script of the same name in \texttt{figures/code/}.

\section*{Figures, tables and equations}

\textbf{Figures}, in the order they appear in the paper.\par
{\centering\footnotesize
\begin{tabularx}{\textwidth}{@{}l l X l@{}}
\toprule
ID & section & content & script (\texttt{figures/code/}) \\
\midrule
F1 & 3.1 & (a) domain, boundary conditions, regions above the plate; (b) hole regions & \texttt{fig\_config\_regions.py} \\
F2 & 3.2 & 13 line classes, plate and centre plane & \texttt{fig\_mesh\_classes.py} \\
F3 & 3.3 & volume fraction above $\deta$, per region, 5 meshes & \texttt{fig\_resolution.py} \\
F4 & 4 & 3D view, RANS 50M against 193M 3DNS: mean $Q$ isosurface, $\eta$ on the wall, cross-planes and insets of $\omega_x$ with streamlines at $x/D$ 2, 4, 10 & \texttt{fig\_overview\_3d.py} \\
F5 & 4.1 & static pressure and reversed fraction along the hole, six rungs & \texttt{fig\_hole\_axis.py} \\
F6 & 4.1 & motion of the diffuser bubble: BR and bubble signal in time, coherence with its 95\,\% level & \texttt{fig\_bubble\_motion.py} \\
F7 & 4.1 & $H$ and $C_f$ along the plate, centreline and mid-pitch & \texttt{fig\_bl\_params.py} \\
F8 & 4.1 & instantaneous 193M centre-plane temperature, $x/D$ $-3$ to 20, with a near-hole view of the shear-layer roll-up & \texttt{fig\_flow\_features.py} \\
F8b & 4.1 & CRVP trajectory and circulation, wall-vortex circulation, RANS 50M against 193M & \texttt{fig\_mean\_features.py} \\
F8c & 4.1 & coolant flux through the top of the near-wall layer $z_x(x)$, mean-flow and turbulent parts, $x/D$ 3--20; RANS/193M ratios of plume half-width, $\bar\eta$ and centreline $\eta$ & \texttt{fig\_flux\_eta.py} \\
F9 & 4.1 & plan view of near-wall $k/U_\infty^2$ and $\gamma$ & \texttt{fig\_wall\_state.py} \\
F10 & 4.1 & transition between the jets: turbulence detector on the $z/D=0.1$ plane, snapshots, $(x,t)$ at mid-pitch and $\gamma(x)$, 193M and 20M LES, with the RANS model $\gamma$ & \texttt{fig\_spots.py} \\
F11 & 4.1 & $\bar\eta$, centreline and mid-pitch $\eta$ against $x/D$, with data & \texttt{fig\_eta\_x.py} \\
F12 & 4.1 & $\eta$ on the wall & \texttt{fig\_eta\_wall.py} \\
F13 & 4.1 & classic $\Pr_t$ of the reference, span-averaged map; grey where not defined & \texttt{fig\_prt\_classic\_span.py} \\
F14 & 4.2 & four stacked maps, $x/D$ 3--20, $z/D$ 0--2: RANS/193M ratio of the conditional span means of $k^{1/2}$, length, $1/\cos\theta$ and $\mu$, with the line $z_x(x)$ & \texttt{fig\_meff\_maps.py} \\
F14b & 4.2 & near-wall layer against $x/D$: exact split of $\ln(\mu_{\mathrm{RANS}}/\mu_{193})$ into density, $k^{1/2}$, $L_{\mathrm{eff}}$ and $\cos\theta$ terms; ratios of $k^{1/2}$, $L$ and $1/\cos\theta$ & \texttt{fig\_meff\_split\_film.py} \\
F15 & 4.2 & classic regional $\Pr_t$ of the 193M against 0.9 and the literature & \texttt{fig\_prt\_regions.py} \\
F16 & 4.2 & $k/U_\infty^2$ and $L_{\mathrm{eff}}/D$ on cross-planes $x/D$ 4, 10, 20 & \texttt{fig\_cross\_planes.py} \\
F17 & 4.2 & length ratio and alignment on the centre plane & \texttt{fig\_centre\_factors.py} \\
F18 & 4.2 & $P_k$, $\varepsilon$ and lengths along wall-normal lines & \texttt{fig\_budget\_lines.py} \\
F19 & 4.2 & spanwise diffusion: directional ratios against the data, lateral flux deficit & \texttt{fig\_prt\_data.py} \\
F20 & 4.3 & $\bar\eta$, centreline, mid-pitch and lateral profile of $\eta$, corrected RANS runs against the 193M & \texttt{fig\_aposteriori.py} \\
\bottomrule
\end{tabularx}\par}

\textbf{Tables}, in the order they appear in the paper.\par
{\centering\footnotesize
\begin{tabularx}{\textwidth}{@{}l l X@{}}
\toprule
ID & section & content \\
\midrule
T1 & 3.1 & rungs and meshes: $s$, cells, solver, model, filter, BR, averaging \\
T2 & 4.3 & a posteriori runs: input changed, recovered deficit, centreline and mid-pitch $\eta$, $\bar\eta$ over $x/D$ 15--20, BR \\
\bottomrule
\end{tabularx}\par}

\textbf{Equations}, in the order they are numbered in the paper.\par
{\centering\footnotesize
\begin{tabular}{@{}l l l@{}}
\toprule
ID & section & quantity \\
\midrule
E1 & 1 & Boussinesq eddy-viscosity relation \\
E2 & 1 & gradient-diffusion hypothesis for the turbulent heat flux, $\Pr_t$ \\
E3 & 3.1 & $\mathrm{Re}_D$, $\theta$ \\
E4 & 3.1 & BR, DR, VR, $I$ \\
E5 & 3.1 & adiabatic effectiveness $\eta$, pitch average $\bar\eta$ \\
E6 & 3.2 & point count per class, $N(s)$ \\
E7 & 3.2 & point placement $\xi_m$ \\
E8 & 3.3 & resolution criterion $\deta\le2$, $\varepsilon$ \\
E9 & 4.1 & $\mu_{\mathrm{eff}}$ and $a^\perp$ \\
E10 & 4.1 & classic $\Pr_t$ and its regional least-squares value \\
E11 & 4.2 & $\mu_{\mathrm{eff}}=\bar\rho k^{1/2}L_{\mathrm{eff}}\cos\theta$ \\
E12 & 4.2 & log split of a ratio into energy, length, alignment \\
E13 & 4.2 & near-wall layer value $\mu_R$ (volume mean over masked points) and its exact four-term split \\
E14 & 4.2 & $\mathrm D\ln k$ budget \\
\bottomrule
\end{tabular}\par}

\textbf{Terminology.} ``Near film'' and ``far film'' for the regions downstream of the hole, never ``wake''.
$\eta$ is the adiabatic effectiveness, $\eta_K$ the Kolmogorov length.

%=====================================================================================================
\section*{1\quad Introduction}
\textbf{Purpose.} Introduction and literature review in one section, no subsections. It narrows from film
cooling in general to the 7-7-7 hole, sets out what is known about RANS and scale-resolving predictions, and
ends on the justification of the configuration of Section 3: the 7-7-7 hole on a flat plate at zero pressure
gradient, with a laminar approach layer, as a test bed for models and meshes at different resolution. The RANS
errors lead the argument. Nine paragraphs, about 2 pages. Prior work in present or past
tense with the author as subject; each paragraph closes on what is missing. The aims follow in Section 2.

\textbf{P1 Film cooling and shaped holes.}
\begin{itemize}
\item HP turbine gas temperatures exceed the material limit; film cooling with compressor bleed air protects the
  surface, and every unit of coolant costs cycle efficiency (Bogard and Thole, 2006).
\item Shaped holes: the diffuser lowers the exit momentum, reduces lift-off and widens coverage, so shaped holes
  give higher $\eta$ than cylindrical holes at high BR (Gritsch et al., 1998a; Thole et al., 1998; Saumweber et
  al., 2003).
\item Cooling design relies on steady RANS because of its cost, and RANS is expected to remain the design tool
  (Milani et al., 2018, 2020); its accuracy for shaped holes therefore decides the design.
\item The accuracy of best-practice RANS for shaped holes is quantified from the literature, as the error in the
  laterally averaged effectiveness $\bar\eta$, so that the RANS error of the Results can be set against it.
\end{itemize}

\textbf{P2 The 7-7-7 baseline hole.}
\begin{itemize}
\item Few hole geometries are shared between research groups, and comparisons against cylindrical holes are of
  limited use (Schroeder and Thole, 2014).
\item Schroeder and Thole (2014) proposed the 7-7-7 as a baseline: $30^\circ$ inclination, $7^\circ$ forward and
  lateral expansion, $L/D=6$, area ratio 2.5, $P/D=6$; $\bar\eta$ peaks near BR 1.5 at DR 1.2 and 1.5.
\item The database since: flow field by PIV and FST effects (Schroeder and Thole, 2016a), thermal field (2017a),
  in-hole roughness (2017b), meter--diffuser offset (Haydt et al., 2017), area ratio (Haydt et al., 2018),
  compound angle (Haydt and Lynch, 2018).
\item Outcome: the most documented shaped hole, with $\eta$, velocity and temperature fields, and the common test
  case for simulations. Justifies the 7-7-7.
\end{itemize}

\textbf{P3 Flat plate at zero pressure gradient.}
\begin{itemize}
\item Generic flat-plate tests leave out acceleration and curvature, so the hole--crossflow interaction is
  studied on its own (Saumweber et al., 2003).
\item All 7-7-7 data are on a flat plate at zero pressure gradient; the 7-7-7 simulations use the same
  configuration (Nunno et al., 2022; Bin and Yang, 2024), as do LES of other shaped holes (Oliver et al., 2017,
  2019; Zamiri et al., 2020).
\item With geometry and pressure gradient fixed, differences between simulations come from the model and the
  resolution alone. Justifies the flat plate at zero pressure gradient.
\end{itemize}

\textbf{P4 RANS predictions and their known errors.}
\begin{itemize}
\item Eddy-viscosity RANS under-predicts the lateral spreading of the coolant and over-predicts the centreline
  $\eta$ (Acharya et al., 2001; Harrison and Bogard, 2008a; Stratton and Shih, 2019); for the 7-7-7, steady RANS
  misses the $\eta$ levels and the offset trends (Haydt et al., 2017).
\item The errors are traced to two closure assumptions, written as E1 and E2.
\item Corrections: anisotropic eddy viscosity (Azzi and Lakehal, 2002; Lakehal, 2002), variable $\Pr_t$ (Liu et
  al., 2008, 2011), higher-order scalar-flux models (Ling et al., 2016a), data-driven closures (Milani et al.,
  2018, 2020; Ellis and Xia, 2024; Laroche et al., 2023).
\item High-fidelity data used to test the assumptions: eddy viscosity and diffusivity are non-uniform and
  anisotropic (Kaszeta and Simon, 2000; Kohli and Bogard, 2005; Coletti et al., 2013); the Boussinesq relation
  does not hold in the hole, along the wall and on the lee side (Ling et al., 2017; Stratton and Shih, 2019).
\item Missing: these studies show where the assumptions fail, not how much of the error each failure carries.
  No study splits the RANS error in the stress and heat flux by region and by cause (turbulence energy, length
  scale, alignment with the mean gradient), or ties each part to the $\bar\eta$ error. The reference for such a
  split has to be resolved; Section 3.3 shows that it is.
\end{itemize}

\textbf{Equations in P4} (standard forms, as in Ling et al., 2016a and Milani et al., 2018).
\begin{equation}
-\bar\rho\,\widetilde{u_i''u_j''}=\mu_t\Big(\frac{\partial\tilde u_i}{\partial x_j}+\frac{\partial\tilde u_j}{\partial x_i}
-\frac{2}{3}\delta_{ij}\frac{\partial\tilde u_k}{\partial x_k}\Big)-\frac{2}{3}\bar\rho k\,\delta_{ij}
\tag{E1}
\end{equation}
Purpose: the Boussinesq relation, one scalar $\mu_t$ aligning the turbulent stress with the mean strain. Names
the assumption that the Results test (isotropy and alignment of the eddy viscosity).
\begin{equation}
-\bar\rho\,\widetilde{u_j''T''}=\frac{\mu_t}{\Pr_t}\,\frac{\partial\tilde T}{\partial x_j},\qquad \Pr_t\approx0.85\text{--}0.9
\tag{E2}
\end{equation}
Purpose: the gradient-diffusion hypothesis with a constant turbulent Prandtl number, which sets the lateral
spreading of the coolant in RANS. The Results test its two inputs: $\mu_t$ (through $\mu_{\mathrm{eff}}$) and the
$\Pr_t$ setting (through the classic $\Pr_t$ of the reference).

\textbf{P5a Mean flow features.}
\begin{itemize}
\item Jets in crossflow carry a counter-rotating vortex pair (CRVP), a horseshoe vortex ahead of the jet and wake
  vortices (Fric and Roshko, 1994; Mahesh, 2013); wake vortices come from the wall boundary layer, so they depend
  on the approach-layer state. In film cooling the attached and detached regimes follow from these structures
  (Bidan et al., 2013; Cavar and Meyer, 2012).
\item Shaped holes: expanded exits reduce the penetration and shear of the jet, with peak turbulence at the exit
  (Thole et al., 1998).
\item Steady RANS holds the mean features only, through the mean strain, and under-predicts the lateral spreading
  they set (Acharya et al., 2001; Harrison and Bogard, 2008a; P4).
\item Missing: how closely the RANS mean features match resolved ones at one hole and one operating point.
  Section 4.1 shows this.
\end{itemize}

\textbf{P5b Instantaneous features.}
\begin{itemize}
\item Jets in crossflow also carry shear-layer roll-up (Kelvin--Helmholtz) and, in film cooling, hairpin vortices
  (Fric and Roshko, 1994; Mahesh, 2013); hairpin packets cause hot spots on the wall (Tyagi and Acharya, 2003).
  Over a laminar approach layer, turbulent spots grow and merge as in bypass transition (Zaki, 2013; Ovchinnikov
  et al., 2008).
\item Whether the structures are coherent depends on the velocity difference across the jet shear layer.
  Fawcett (2011) found shear-layer and hairpin vortices downstream of cylindrical and shaped holes only when that
  difference was large; at BR 0.8 and 1.0 (DR 1, so VR near 1) the shaped hole showed no coherent unsteadiness
  in the jet. With turbulent crossflow the hairpin peak flattens to a weak pressure peak (Ziefle and Kleiser,
  2013).
\item The present case has $\mathrm{VR}=\mathrm{BR}/\mathrm{DR}\approx1.0$, so the literature predicts a jet
  with weak coherence; the mixing is then mostly broadband turbulence, which is the part E1 and E2 are written
  for.
\item Missing: the structures are known; how much of the mixing each mechanism carries, and how a model
  represents that share, is not. The question for this paper is the share of the physics attributed to each
  mechanism, not the share of the unsteadiness a mesh resolves.
\end{itemize}

\textbf{P6 Scale-resolving simulations and their resolution.} One paragraph, no table: the studies grouped by
resolution and modelling, then what each group found compared and contrasted.
\begin{itemize}
\item Coarse LES and hybrid RANS, $1$--$3\times10^{6}$ cells: Tyagi and Acharya (2003), Ziefle and Kleiser
  (2013), Foroutan and Yavuzkurt (2015), Yu and Yavuzkurt (2018). Found: hairpin packets set the hot spots; the
  crossflow turbulence decides the spreading, and without it $\eta$ is wrong; DES mixes more than RANS
  wall-normally and laterally; URANS centreline $\eta$ is off by $\sim$50\,\% against $\sim$11\,\% for LES.
  Resolution: grid studies on the RANS, or none; Ziefle and Kleiser call their grid coarse.
\item Wall-resolved LES of round holes with turbulent inflow, $4$--$10\times10^{6}$ points: Iourokina and Lele
  (2006), Peet and Lele (2008), Johnson and Kapat (2013). Found: in-hole separation and plenum effects; the
  isotropic $\mu_t$ and $\alpha_t$ do not hold (Peet and Lele); the errors left are near the injection and the
  wall. Resolution judged in wall units or by the resolved fraction of $k$.
\item LES of shaped holes, $6$--$54\times10^{6}$ points: Zamiri et al.\ (2020), Oliver et al.\ (2017, 2019), Nunno
  et al.\ (2022), the last three on the 7-7-7 at BR 2 with a laminar approach layer. Found: pulsation inside the
  hole; backflow at the exit; the gradient-diffusion heat flux is wrong (Oliver et al., 2017); Mach number changes
  $\bar\eta$ by up to 50\,\%. Oliver et al.\ (2017) put their near-hole error ($\Delta\eta\approx0.05$) down to the
  resolution. Resolution in wall units ($\Delta x^+\approx12$--34).
\item DNS, $78\times10^{6}$ cells: Bin and Yang (2024), the 7-7-7 across compound angles, turbulent inflow at
  BR 1. Resolution in wall units and $\Delta<5\eta_K$ outside the wall layer; no closure analysis.
\item Contrast: each group finds that resolving more of the flow improves $\eta$, but the groups differ in
  model, mesh and resolution measure at once, so the gain cannot be assigned to either. A closure is assessed in
  three studies (Peet and Lele, 2008; Oliver et al., 2017; Foroutan and Yavuzkurt, 2015), none against a resolved
  DNS.
\item Missing: no study takes one shaped hole through RANS, LES and DNS on one mesh family, so the resolution
  needed to predict $\bar\eta$, and the share of the RANS error each closure assumption carries, are not known.
\end{itemize}

\textbf{P7 The laminar approach boundary layer.}
\begin{itemize}
\item Most studies, and the 7-7-7 experiment, use a tripped turbulent approach layer.
\item On the suction surface of cooled HP rotors the boundary layer is transitional for
  $3\times10^5<\mathrm{Re}_x<10^6$, and at the lower $\mathrm{Re}_x$ of small cores the layer at a cooling row
  can be laminar (Clark et al., 2024); whether a jet trips it depends on the size, spacing and position of the
  perturbation (Kar et al., 2026). These two references justify the laminar approach layer; no figure is
  needed.
\item Laminar-approach film-cooling data are rare: Goldstein and Yoshida (1982) found that the holes trip the
  laminar layer, and that the injectant state matters more than the approach state.
\item Freestream turbulence drives bypass transition of the approach layer (Jacobs and Durbin, 2001; Brandt et
  al., 2004); RANS needs a transition model to hold the layer laminar (Langtry and Menter, 2009).
\item Justifies the laminar approach layer with FST, and the $\gamma$--$\mathrm{Re}_\theta$ model in RANS. It is
  also the condition furthest from the turbulent flows the closures were calibrated on.
\end{itemize}

\textbf{P8 The present study.}
\begin{itemize}
\item The justification chain in one place: the RANS errors of P4 are known by place but not by share (P4);
  at VR near 1 the mixing is mostly broadband turbulence, the part E1 and E2 are written for, so the question is
  how much of the mixing each mechanism carries (P5); scale-resolving simulations give the reference, at a
  resolution not yet established (P6); the 7-7-7 is the most documented
  shaped hole (P2); a flat plate at zero pressure gradient fixes geometry and pressure gradient (P3); a laminar
  approach layer is relevant to small cores and hardest for RANS (P7). Together they make a test bed in which
  only the model and the mesh change.
\item The approach: one mesh family from 5M to 193M cells; steady RANS, LES and DNS; the same boundary conditions
  and plenum pressure on every rung. BR is then an output of each rung, and it differs between rungs through the
  separation in the diffuser of the hole (Haydt et al., 2017, 2018), which 4.1 introduces as a short lead-in.
\item Closing sentence leads into Section 2, the research aims.
\end{itemize}

%=====================================================================================================
\section*{2\quad Research Aims}
\textbf{Purpose.} A short section between the Introduction and the Methods, as in the GPPS paper: one framing
sentence (``This paper aims to \dots''), then the aims as a numbered list, one sentence each. No figures, tables
or equations.
\begin{enumerate}[label=\arabic*.]
\item Quantify the $\bar\eta$ error of each rung, RANS to 100M 3DNS, against the 193M reference.
\item Attribute the RANS error to the closure inputs: the eddy-viscosity factors ($\mu_t$), the $\Pr_t$ setting
  and the spanwise diffusion (E1, E2).
\item Test a posteriori which correction recovers $\bar\eta$: the $\Pr_t$ setting, QCR, heat-flux substitution and
  $\gamma=1$.
\item State the resolution each mesh achieves and the resolution needed to predict $\bar\eta$ (supporting aim).
\end{enumerate}
The conclusions mirror the aims.

%=====================================================================================================
\setcounter{figure}{0}
\setcounter{table}{0}
\section*{3\quad Methods}

%-----------------------------------------------------------------------------------------------------
\subsection*{3.1\quad Simulation Configuration and Computational Methods}
\textbf{Purpose.} Define what is simulated, with which codes and models, under which boundary conditions, and
how the domain is divided for analysis. After this subsection the reader can place any rung name or region name
used later in the paper.

\textbf{(i) Geometry and operating point.}
\begin{itemize}
\item Baseline 7-7-7 laidback fan-shaped hole (Schroeder and Thole, 2014): $D=7.75$\,mm, $30^\circ$
  inclination, pitch $6D$; single hole, periodic in $y$; coolant fed at the metering inlet face, no plenum.
\item Flow conditions, justified by the transitional HP turbine (Clark et al., 2024; Kar et al., 2026): Blasius
  layer, $\mathrm{Re}_\theta=180$ at the hole, $M_\infty=0.25$, Tu 5.5\,\% with $L=0.43D$, $\mathrm{Re}_D=2408$;
  raised viscosity $\mu_\infty=2.5\times10^{-4}$\,Pa\,s; adiabatic walls.
\item BR 1.5 at DR 1.5: peak of $\bar\eta$ in the 7-7-7 database. Definitions of BR, DR, VR, $I$, $\eta$,
  $\bar\eta$, $D/U_\infty$, FT. All rungs share the plenum $p_{0,c}$, so BR is an output of each rung (T1).
\end{itemize}

\textbf{Equations in (i).}
\begin{equation}
\mathrm{Re}_D=\frac{\rho_\infty U_\infty D}{\mu_\infty}=\frac{0.9003\times86.28\times0.00775}{2.5\times10^{-4}}=2408,
\qquad \theta=0.075D
\tag{E3}
\end{equation}
Purpose: fixes the Reynolds number of the hole and the approach layer, and justifies the raised viscosity
($\sim$14 times that of air) needed to reach it at the physical hole size and $M_\infty=0.25$.
\begin{equation}
\mathrm{BR}=\frac{\dot m_c/A_m}{\rho_\infty U_\infty},\qquad A_m=\frac{\pi D^2}{4},\qquad
\mathrm{DR}=\frac{\rho_c}{\rho_\infty},\qquad
\mathrm{VR}=\frac{\mathrm{BR}}{\mathrm{DR}},\qquad I=\frac{\mathrm{BR}^2}{\mathrm{DR}}
\tag{E4}
\end{equation}
Purpose: defines the operating point as the experiment does, so BR 1.5, DR 1.5 compares directly with the
7-7-7 database; $A_m$ is the metering cross-section.
\begin{equation}
\eta=\frac{T_{0,\infty}-T_w}{T_{0,\infty}-T_{0,c}},\qquad
\bar\eta(x)=\frac{1}{P}\int_{-P/2}^{P/2}\eta\,\mathrm{d}y,\qquad P=6D
\tag{E5}
\end{equation}
Purpose: adiabatic effectiveness and its pitch average, the performance measure used throughout the Results.
$\eta$ is reserved for the effectiveness; the Kolmogorov length is $\eta_K$ (E8).

Not given as equations: $D/U_\infty=89.8$\,\textmu s and $1\,\mathrm{FT}=40D/U_\infty=3.59$\,ms (inline), and
Sutherland's law in 3DNS (named and cited, no formula).

\textbf{(ii) Computational methods.}
\begin{itemize}
\item One mesh family, one set of boundary conditions, one reference (193M 3DNS). RANS against 50M 3DNS on the
  same mesh isolates modelling; 50M, 100M and 193M 3DNS isolate resolution; the 20M LES differs in both (T1).
\item STAR-CCM+ RANS: coupled density-based solver, SST $k$--$\omega$ with $\gamma$--$\mathrm{Re}_\theta$, which
  keeps the approach layer laminar ($\mathrm{Re}_\theta=176$, $H=2.51$--2.64 at $x/D=-4$); max $y^+=1.1$;
  residuals down 6--8 decades. RANS on the 5M and 50M meshes agree to 0.06\,\% in BR. URANS returns the
  steady state and is not a rung.
\item Every rung, including the RANS, runs at the same plenum $p_{0,c}$. BR matching is not mentioned in the
  Methods: the need for it is introduced in the Results, where the differences in diffuser-bubble behaviour
  between rungs are described.
\item 3DNS: high-order compressible finite-difference code (Wheeler et al.); 7-point DRP stencil (1\,\% phase
  error at $\sim$4 points per wavelength; Tam and Webb, 1993); 8th-order filter, $sf=0.05$ on 50M (implicit
  LES), 0.03 on 100M and 193M; 20M LES with WALE.
\item Initialisation, flush (BR drift below 1\,\% over 0.5\,FT) and averaging time.
\end{itemize}

\textbf{Table T1: rungs and meshes.} One row per mesh of the family, with its scale factor and the
rungs run on it. Purpose: the single reference for every rung name and every mesh; $s$ links it to the
coarsening in 3.2. No machine or rank information.

\begin{table}[H]
\centering\footnotesize\setlength{\tabcolsep}{4pt}
\begin{tabular}{@{}l c c l l l c c@{}}
\toprule
mesh & $s$ & cells & solver & model & filter $sf$ & BR & averaging \\
\midrule
5M   & 0.30$^a$ & $5.77\times10^{6}$ & STAR-CCM+ & RANS, SST + $\gamma$--Re$_\theta$ & -- & 1.591 & steady \\
20M  & 0.453 & $2.31\times10^{7}$ & 3DNS & LES, WALE & 0.03 & 1.498 & 6\,FT \\
50M  & 0.628 & $5.31\times10^{7}$ & STAR-CCM+ & RANS, SST + $\gamma$--Re$_\theta$ & -- & 1.592 & steady \\
     &       &                    & 3DNS & none (implicit LES) & 0.05 & 1.516 & 6\,FT \\
100M & 0.796 & $1.13\times10^{8}$ & 3DNS & none & 0.03 & 1.526 & 6\,FT \\
193M & 1     & $1.92\times10^{8}$ & 3DNS & none (reference) & 0.03 & 1.521 & 10\,FT \\
\bottomrule
\end{tabular}
\caption{Rungs and meshes. $s$: scale factor of E6 (3.2). Cells: summed over the 23 blocks of each mesh. BR:
coolant mass flow through the diffuser sections at the common plenum $p_{0,c}$. Averaging: 193M batch b010.
$^a$\,The 5M mesh has the blocks of the family, but its counts were not set by E6; 0.30 is the average of
$(N_{5}-1)/(N_{193}-1)$ over the blocks, weighted by the 193M points in each block (per-line range 0.20--0.46).}
\end{table}

\textbf{(iii) Boundary conditions.} The same for every rung, so stated once, in text only (no table, no
separate figure; F1(a) marks where each applies):
\begin{itemize}
\item Inlet: Blasius layer, $\mathrm{Re}_\theta$ 135--140, with FST of Tu 5.5\,\%; outlet and top at
  $p_\infty$; adiabatic no-slip walls; periodic span; coolant inlet at the same $p_{0,c}$ and $T_{0,c}$ for
  every rung, with coolant turbulence.
\item How each code imposes them, in one sentence each: FST box and characteristic boundaries with buffers in
  3DNS, inlet intensity in STAR-CCM+; Sutherland viscosity in 3DNS, constant in STAR-CCM+.
\end{itemize}

\textbf{(iv) Analysis regions.} Eight regions that fill the domain without overlap (buffers excluded):
metering, diffuser (split $3.5D$ along the axis from the breakout), exit (breakout face extruded $0.5D$ each
way), approach layer, jet ($x/D$ $-1.00$ to 2.36), near film (to 10.25), far film (to 20), freestream
($z/D\ge4$). One sentence per boundary choice. Used in 3.3 and in every regional average of the Results.

\textbf{Figure F1: where the domain, boundary conditions and regions are.} Configuration of the domain and the
hole, full page width.
\begin{itemize}
\item WHY: one picture for (i), (iii) and (iv), so that every rung name, boundary condition and region name used
  later can be placed.
\item HOW: panel (a) gives the domain extent, the hole position, the condition on every boundary and the five
  regions above the plate. Panel (b) enlarges the hole to show the metering, diffuser and exit regions and the
  coolant inlet face. The 3DNS buffers are hatched.
\item WHAT: every region is named on the figure, so colour is not the only key (Okabe-Ito fills). The hatched
  buffers show which parts of the domain are excluded from all data.
\end{itemize}
\begin{figure}[H]
\centering
\includegraphics[width=0.92\textwidth]{fig_config_regions.pdf}
\caption{F1 (\texttt{figures/code/fig\_config\_regions.py}). (a) Centre plane of the domain, to scale,
with the boundary conditions and the regions above the plate. (b) Enlarged view of the hole with the hole
regions. Hatched: 3DNS buffer zones.}
\end{figure}

%-----------------------------------------------------------------------------------------------------
\subsection*{3.2\quad Mesh Family and Coarsening Method}
\textbf{Purpose.} Show that the meshes differ in resolution only: same blocks, walls and hole geometry,
obtained by coarsening the 193M mesh line by line.

\textbf{Content.}
\begin{itemize}
\item The 193M mesh: 23 blocks; 80--100\,\um\ across the jet; 220--290\,\um\ axially in the hole; 40\,\um\ first
  cell; 1\,\% growth to $z/D=10$.
\item Coarsening: 13 classes of connected grid lines, one count per class; E6 with the $s$ of T1;
  points evenly spaced in 193M point number (E7), so the 193M stretching is kept and each cell grows by
  $\sim1/s$; wall cell kept; interface jumps below $\sim$1.5; cubic placement on the 193M line.
\item Departures from E6: classes 1, 9, 11 held at the 100M count; 2, 7, 8, 10 raised for load balance.
\item The 5M mesh shares the 23 blocks but was not built with E6; its representative $s$ is given in T1 and the
  range of its per-line ratios in the text.
\end{itemize}

\textbf{Equations.}
\begin{equation}
N=\operatorname{round}\!\big((N_{193}-1)\,s\big)+1,\qquad \text{cells}\approx s^3\times1.92\times10^{8}
\tag{E6}
\end{equation}
Purpose: the point count of every class on a coarser mesh from its 193M count $N_{193}$ and the one scale
factor $s$ of that mesh (T1); the second part links $s$ to the mesh size.
\begin{equation}
\xi_m=m\,\frac{N_{193}-1}{N-1},\qquad m=0,\dots,N-1
\tag{E7}
\end{equation}
Purpose: where the $N$ points of a coarser line sit, as fractional 193M point numbers $\xi$; even spacing in
$\xi$ keeps the 193M stretching and cell aspect ratios. The wall-cell and interface exceptions are stated in the
text, not as equations.

\textbf{Figure F2: how the 13 line classes sit in the mesh.} The 13 classes of grid lines on the plate plane and
the centre plane.
\begin{itemize}
\item WHY: makes ``class'' concrete; E6 and the class numbers quoted in the text refer to it.
\item HOW: (a) the plate plane with all 23 blocks; (b) the centre plane through the hole, both of the 20M mesh.
  Each class is drawn in an Okabe-Ito colour, with a line style per grid direction (solid $i$, dashed $j$,
  dash-dot $k$).
\item WHAT: the figure shows which lines belong to which class, so the classes that depart from E6 (1, 9, 11, 2,
  7, 8, 10) can be named by number in the text.
\end{itemize}
\begin{figure}[H]
\centering
\includegraphics[width=0.85\textwidth]{fig_mesh_classes.pdf}
\caption{F2 (\texttt{figures/code/fig\_mesh\_classes.py}, 20M mesh). (a) Plate plane, all 23 blocks.
(b) Centre plane through the hole.}
\end{figure}

\textbf{In text, no table:}
\begin{itemize}
\item the classes that depart from E6 and why, by class number on F2: 1, 9, 11 held at the 100M counts across
  the hole (88, 215, 48 points); 2, 7, 8, 10 raised for load balance;
\item the 5M mesh: same blocks, per-line ratios $(N_5-1)/(N_{193}-1)$ of 0.20--0.46;
\item the cell sizes that stay fixed or bounded: first cell 40\,\um\ at the plate and 22\,\um\ in the hole on
  every mesh, cell-size jump across block interfaces below $\sim$1.5;
\item the 193M cell sizes (80--100\,\um\ across the jet, 220--290\,\um\ axially in the hole, 0.59\,mm at the top),
  which grow by $\sim1/s$ on the coarser meshes.
\end{itemize}

%-----------------------------------------------------------------------------------------------------
\subsection*{3.3\quad Achieved Mesh Resolution}
\textbf{Purpose.} State, for each mesh and region, how the grid size compares with the smallest scale of the
reference flow.

\textbf{Equation.}
\begin{equation}
\frac{\Delta_{\mathrm{mesh}}}{\eta_{K,193}}\le2,\qquad \Delta_{\mathrm{mesh}}=|J|^{1/3},\qquad
\eta_{K,193}=\Big(\frac{\nu^3}{\varepsilon/\bar\rho}\Big)^{1/4},\qquad
\varepsilon=\overline{\tau_{ij}\partial_ju_i}-\bar\tau_{ij}\,\partial_j\bar u_i
\tag{E8}
\end{equation}
with $|J|$ the cell volume of the 5M, 20M, 50M, 100M or 193M mesh, $\nu=\bar\mu/\bar\rho$, and $\varepsilon$ the
viscous dissipation of the fluctuations, all from the 193M 3DNS, statistics batch b010 (10\,FT).
Purpose: the resolution criterion and the one quantity reported in F3 and the text. The bound 2 follows from
$\kappa_{\max}\eta_K\ge1.5$ for spectral DNS (Pope, 2000), i.e.\ $\Delta/\eta_K\le\pi/1.5\approx2.1$; stated
inline, not numbered.
\begin{itemize}
\item One $\eta_K$ field for all meshes. The ratio measures each mesh against the scales of the reference flow,
  so differences between meshes come from the mesh alone. The 5M and 50M RANS meshes are rated on the same
  scale as the scale-resolving rungs.
\item $\eta_{K,193}$ is computed on the 193M mesh, so where the 193M is not resolved (diffuser) it is
  less certain.
\item Wall units: one sentence giving the maximum $y^+$, $x^+$ and $z^+$ on each mesh.
\item Convergence of the statistics between 100M and 193M is a separate criterion, referenced here.
\end{itemize}

\textbf{Method to compute it.} $\Delta_{\mathrm{mesh}}$ of each coarser mesh is mapped
onto the 193M points through the line-by-line correspondence of the family, and volume fractions are taken with
193M cell volumes, so every mesh uses the same region volumes. The 5M mesh is mapped through its own per-line
counts.

\textbf{In text, no table:}
\begin{itemize}
\item volume fraction with $\deta>2$ over the whole domain, per mesh;
\item the regions that are not resolved and by how much: diffuser, exit and far film on each mesh; the metering
  section in one sentence, since it is left out of F3;
\item $\eta_{K,193}$ in the diffuser ($\sim$47\,\um) against its 220--290\,\um\ axial cells, as the reason the
  diffuser is the least resolved region.
\end{itemize}

\textbf{Figure F3: how far each mesh is from the resolution criterion.} Fraction of each region's volume above a
given $\deta$, one line per mesh: 5M, 20M, 50M, 100M, 193M.
\begin{itemize}
\item WHY: the main evidence of 3.3. The figure shows how far each mesh of the family sits from the criterion of
  E8, region by region, from the RANS mesh to the reference, including the tails behind the values quoted in the
  text.
\item HOW: $\deta$ (E8) is evaluated at every cell of the 193M mesh for each of the five meshes, and the volume
  fraction above each value is plotted. One row of panels: diffuser, exit, near film, far film and the whole
  domain. The metering section is left out and covered in one sentence of the text. The five meshes are
  distinguished by Okabe-Ito colours, line styles (solid, dashed, dash-dot, dotted) and one marker style.
\item WHAT: the diffuser is the least resolved region on every mesh, since $\eta_{K,193}$ there ($\sim$47\,\um) is
  small against its 220--290\,\um\ axial cells. A dashed vertical line marks the ratio 2.
\end{itemize}
\begin{figure}[H]
\centering
\includegraphics[width=\textwidth]{fig_resolution.pdf}
\caption{F3 (\texttt{figures/code/fig\_resolution.py}): volume fraction above $\deta$ in the diffuser, exit, near
film, far film and whole domain; 20M LES b006, 50M 3DNS b006, 100M 3DNS b006 (6\,FT each), 193M 3DNS b010 (10\,FT).
Dashed vertical line: ratio 2.}
\end{figure}

%=====================================================================================================
\section*{4\quad Results}
\textbf{Purpose.} The Results move from what the reader can see to what the closure does. Section 4.1 gives
the coolant flow in detail: the hole, the mean flow features of the jet in the RANS and in the reference, the
mainstream layer between the jets, and the effectiveness they produce. Section 4.2 breaks the mean mixing of each rung down into its
causes and compares them with the literature. Section 4.3 gives the a posteriori tests.

\textbf{Common to 4.1 and 4.2.}
\begin{itemize}
\item Rungs: the six of Section 3 only: RANS 5M and 50M, 20M LES, 50M, 100M and 193M 3DNS, all at the same
  plenum $p_{0,c}$. No corrected RANS (QCR, $\Pr_t$ sweep, substitution) before 4.3.
\item Reference: the 193M 3DNS at batch b010 (10\,FT), the record it was run for, in every figure and number.
  The lower rungs keep their last batch (b006, 6\,FT); the record-length check (193M b006 against b010: regional
  values move by 0.046 or less) is one sentence.
\item Quantities the reader knows are shown beside the new ones: $k/U_\infty^2$, shape factor $H$, $C_f$,
  $\mathrm{Re}_\theta$, intermittency $\gamma$, $k$ production $P_k$ and dissipation $\varepsilon$ (normalised by
  $U_\infty^3/D$), the dissipation length $k^{3/2}/\varepsilon$, the discharge coefficient $C_D$ and the static
  pressure along the hole. All 3DNS values come from the time-averaged statistics (means, Favre stresses, heat
  fluxes, $\widetilde{T''^2}$, dissipation, pressure--dilatation, scalar dissipation); RANS values from the cell
  exports.
\item Line plots in PDF, fields in PNG; $\eta$ and temperature fields in \texttt{jet}, other sequential fields in
  \texttt{inferno}, signed fields and log ratios in \texttt{bwr} with symmetric limits.
\item Every comparison with a reference is against the 193M 3DNS, b010.
\end{itemize}

\textbf{Figure F4: the flow in three dimensions, RANS against the reference.} First figure of Section 4: RANS 50M
and 193M 3DNS in one view, same camera.
\begin{itemize}
\item WHY: orients the reader at a glance and shows the differences in the mean flow before any section analyses
  them.
\item HOW: (a) RANS 50M, (b) 193M 3DNS (b010), same camera and same threshold. Mean-flow isosurface
  $Q(D/U_\infty)^2=0.005$ in uniform grey (opacity 0.8); $\eta$ (E5) on the wall in \texttt{jet}; the hole
  outline; cross-planes at $x/D$ 2, 4, 10 with $\omega_xD/U_\infty$ (\texttt{bwr}, $\pm0.15$) and in-plane
  streamlines, and the frames of the centre plane and of $x/D=20$. Each cross-plane is repeated as a flat inset
  ($|y/D|\le3$, $z/D$ 0--2) joined to its plane by a line. Insets are viewed looking downstream, as the 3D view. Labels by acronym: CRVP and HV in the view, WV in
  the insets. The 193M mean vorticity on the planes is smoothed over 0.06\,$D$. Format PNG.
\item WHAT: where the CRVP and the vortex around the upstream edge of the hole (HV) lie in each rung; the
  wall vortices (WV) at $x/D$ 2, 4, 10; how the footprint of $\eta$ differs between RANS and reference.
  The isosurfaces reach $x/D$ 13.2 (RANS) and 14.7 (193M).
\end{itemize}
\planfig{fig_overview_3d.png}{0.8}{F4 (\texttt{figures/code/fig\_overview\_3d.py}): (a) RANS 50M, (b) 193M 3DNS
(b010), same camera; isosurface $Q(D/U_\infty)^2=0.005$, $\eta$ on the wall, hole outline; cross-planes and
insets of $\omega_xD/U_\infty$ with in-plane streamlines at $x/D$ 2, 4, 10.}

%-----------------------------------------------------------------------------------------------------
\subsection*{4.1\quad Coolant Flow}
\textbf{Purpose.} The coolant flow in detail, and the comparison of the mean flow of the RANS with that of the
193M 3DNS. Four steps: the hole (a short lead-in that explains the BR difference between rungs), the mean flow
of the jet above the wall (counter-rotating vortex pair, secondary vortices, upwash, jet trajectory), the
mainstream layer between the jets, and the effectiveness that these produce. The section ends with the frame in
which the rest of the paper attributes the differences. All differences in the mean flow features between the
RANS and the reference are analysed here; 4.2 takes the closure only. About 3 pages.

\textbf{(i) First chain, a short lead-in: from the plenum face through the hole.}
\begin{itemize}
\item The coolant enters at the plenum face at the same $p_{0,c}$ (81\,110\,Pa) on every rung, so BR is a
  result. It accelerates through the metering section and separates in the diffuser.
\item The diffuser separation depends on the modelling: the 3DNS hold two corner lobes on the downstream wall
  that grow from $s/D=-3.3$, peak at 16--19\,\% of the section near $s/D=-2.1$ and close at the exit; the RANS
  keeps two thin lobes of 9--10\,\%; the 20M LES lobes peak at 21\,\%.
\item The reversed area fraction on the sections at $s/D$ $-2.3$, $-1.1$, $-0.5$ is 0.092, 0.089, 0.051 (RANS
  50M), 0.179, 0.148, 0.027 (20M LES) and 0.159, 0.100, 0.028 (193M); the centre-plane reversed layer is
  0.03--0.07\,$D$ thick, with maximum thickness 0.032\,$D$ (RANS), 0.072\,$D$ (20M LES) and 0.048\,$D$ (193M).
\item The separation sets the pressure recovery of the diffuser, and through it the head across the metering
  section and the flow. The static rise from the metering section to the exit is 2627--2645\,Pa in the RANS
  against 1148--1554\,Pa in the scale-resolving rungs; the RANS also loses about 600\,Pa more in the metering
  section.
\item Result: BR 1.591--1.592 (RANS 5M and 50M, mesh-independent, so a model effect), 1.498 (20M LES), 1.516,
  1.526, 1.521 (50M, 100M, 193M 3DNS); the RANS passes 4.7\,\% more coolant, the 20M LES 1.5\,\% less. $C_D$ per
  rung in the text.
\item Literature: the diffuser recovery and the conical-diffuser analogy (Haydt et al., 2018); $C_D$ of
  expanded holes (Gritsch et al., 1998b); backflow at the 7-7-7 exit in iLES (Oliver et al., 2017). Not seen
  before: the rung-to-rung BR spread at one plenum pressure, traced to the diffuser recovery.
\item Bubble motion (193M b010, F6): the separated volume is $0.719\,D^3$ and varies by 10\,\% (standard deviation); the two lobes
  trade volume (correlation $-0.48$); BR is $1.527\pm0.019$ (1.2\,\%). At St 0.273 BR is coherent with the total
  volume (0.86, 95\,\% level 0.41) and lags it by 43--46$^\circ$; it is not coherent with the lobe difference
  (0.08). BR follows the size of the bubble, not the exchange between the lobes.
\item Effect on the mainstream (193M b010): the share of the energy of the slow bubble mode that sits in the
  film, and the coherence of the lobe switching with the film meander and with the sideways shift of the cooled
  footprint. The expectation of Fawcett (2011) for VR near 1 is that the bubble motion does not organise the
  film.
\item BR sensitivity in one sentence from the 7-7-7 data at DR 1.5: $\bar\eta$ ($x/D$ 7--15) 0.230, 0.243, 0.218
  at BR 1.0, 1.5, 2.0, so a 4.7\,\% BR difference moves $\bar\eta$ by much less than the RANS deficit of (iv).
\end{itemize}

\textbf{Figure F5: where the rungs part along the hole.} Pressure and separation against $s/D$, six rungs.
\begin{itemize}
\item WHY: BR differs between rungs at one plenum pressure. The figure shows where along the hole the rungs part
  and what sets the difference.
\item HOW: both panels are plotted against $s/D$ for all six rungs, with the metering section and the diffuser
  marked. (a) The area-mean static pressure coefficient $C_{p,c}=(p-p_\infty)/(p_{0,c}-p_\infty)$. (b) The
  fraction of each section with reversed axial flow.
\item WHAT: the scale-resolving rungs share the same metering drop, the RANS has a larger one, and the diffuser
  recovery follows the size of the reversed region. The static rise is 2627--2645\,Pa in the RANS against
  1148--1554\,Pa in the scale-resolving rungs. For the 193M: BR 1.520, peak reversed fraction 0.159 at
  $s/D=-2.3$, diffuser static rise 1537\,Pa ($s/D$ $-4.9$ to $-0.9$); b006 is within 0.1\,\% in BR and 7\,Pa in
  the rise.
\end{itemize}
\planfig{fig_hole_axis.pdf}{0.5}{F5 (\texttt{figures/code/fig\_hole\_axis.py}): (a) $C_{p,c}$, (b) reversed-flow
fraction of each section, against $s/D$; all six rungs, 193M b010, other 3DNS and LES b006.}

\textbf{Figure F6: how the diffuser bubble moves and whether BR follows it.} 193M 3DNS, b010: BR and bubble signal
in time, and their coherence.
\begin{itemize}
\item WHY: shows that the bubble moves and how closely the BR fluctuation follows it, so that the BR of each
  rung can be attributed to the mean separation and not to its unsteadiness.
\item HOW: both panels use the long continuous stretch of the 193M record. (a) BR and the bubble signal on the
  same time axis (103.8\,$D/U_\infty$, 519 records). The bubble signal is the total separated volume $V$; the
  volume of each lobe ($V_+$, $V_-$: the reversed-flow regions in the two lower corners of the diffuser, $y>0$
  and $y<0$) and the lobe difference $V_+-V_-$ are given in the text. (b) One line, the magnitude-squared coherence $\gamma^2(f)$ between BR and the bubble
  signal, between 0 and 1: the share of the BR fluctuation at frequency $f$ that is linearly related to the
  bubble signal at the same frequency. The 95\,\% level is drawn as one clear line; it is the value two
  unrelated signals would exceed by chance in only 5\,\% of cases, given the number of averaged segments
  ($1-0.05^{1/(n-1)}$ for $n$ independent segments, with the reference of the method). The frequencies where the
  coherence exceeds it are marked. The St axis is cut at 1, since above it the spectra carry no variance of note.
\item WHAT: over the stretch BR is $1.527\pm0.019$ and $V$ is $0.729\pm0.064\,D^3$. The coherence of BR with $V$
  exceeds the 95\,\% level (0.41, 7 segments) in four bins, St 0.195--0.313, and peaks at 0.86 at St 0.273. With
  the single lobes it is 0.87 and 0.71; with the lobe difference 0.08. BR follows the size of the bubble near
  St 0.27 and varies by 1.2\,\%.
\end{itemize}
\planfig{fig_bubble_motion.pdf}{0.7}{F6 (\texttt{figures/code/fig\_bubble\_motion.py}): (a) BR and bubble signal
against time, (b) coherence between them with the 95\,\% level (line) and the frequencies above it (marks);
St $\le1$. 193M 3DNS, b010, long continuous stretch.}

\textbf{(ii) The jet above the wall: mean flow features, RANS against the reference.}
\begin{itemize}
\item Names used from here on: counter-rotating vortex pair (CRVP); horseshoe vortex (HV), the vortex around the
  upstream edge of the hole exit seen in the isosurface of F4 (as named for this hole by Tracy and Lynch, 2020);
  wall vortex (WV), the near-wall vortex outboard of the CRVP with the opposite sign to the CRVP vortex on its
  side. WV is a descriptive name defined here; the anti-kidney pair of Haven and Kurosaka (1997) has the same
  sign but lies above the jet, and is noted as a possibly related structure. Each is defined once in the text
  against F4. Cross-planes are viewed looking downstream throughout.
\item Origin of the CRVP: the vorticity of the hole boundary layers and of the diffuser side walls leaves the exit
  as two streamwise vortices, which carry hot gas down at the flanks and coolant up on the centreline. One
  sentence, with Haven and Kurosaka (1997) and Fric and Roshko (1994) for the mechanism and Haydt et al. (2018)
  for the shaped hole.
\item CRVP trajectory: lateral position $y_c(x)$ and height $z_c(x)$ of each core from $x/D$ 2 to 20, RANS 50M
  against 193M (vorticity-weighted centroid inside half the peak, on the fine cross-planes, same smoothing of
  0.15\,$D$ in both rungs). The RANS pair starts wider apart ($|y_c|/D$ 0.75 against 0.60 at $x/D=3$) and
  agrees within 0.05\,$D$ from $x/D=6$. The RANS cores are lower throughout: $z_c/D$ 0.40, 0.54, 0.68 against
  0.48, 0.67, 0.80 at $x/D$ 4, 10, 20.
\item CRVP strength: circulation $\Gamma(x)/(U_\infty D)$ of each vortex. The RANS circulation is 0.166, 0.141,
  0.085, 0.048 against 0.093, 0.096, 0.066, 0.041 at $x/D$ 3, 4, 10, 20: 1.8 times the reference at $x/D=3$,
  1.3 times at 10 and 1.16 times at 20. The RANS pair starts stronger and decays faster; the reference
  circulation is flat to $x/D=5$. The decay of $\Gamma$ is compared with the turbulent diffusion each rung has
  (link to $\mu_{\mathrm{eff}}$ of 4.2).
\item Secondary vortices, from the fine cross-planes at ten stations with a smoothing of 0.02\,$D$ (they are
  0.1--0.3\,$D$ thick): the 193M WV holds a circulation of 0.015--0.025 from $x/D$ 3 to 15 while its peak falls
  from 0.50 to 0.17; at $x/D=20$ it is no longer a closed vortex. The RANS WV has a circulation of 0.029 and
  0.021 at $x/D$ 3 and 4 (peak 0.48 at $x/D=4$) and is gone from $x/D=5$, where the vorticity is a wall-attached
  layer. A weaker vortex of the same sign as the CRVP lies further outboard in the RANS at $x/D=4$; in the 193M
  it is at the noise floor of the mean field and is not named or quantified. The asymmetry of the 193M is stated
  with the record length.
\item Upwash and the coolant leaving the near-wall layer. Coolant scalar from the total enthalpy,
  $c=(h_{0,\infty}-h_0)/(h_{0,\infty}-h_{0,c})$. The time-averaged vertical flux splits (Favre) as
  $\langle\rho wc\rangle=\langle\rho w\rangle\tilde c+\langle\rho w''c''\rangle$: the first term is the mean-flow
  upwash, the second the turbulent flux, measured in the 193M and modelled in the RANS as
  $-(\mu_t/\Pr_t)\,\partial\tilde c/\partial z$ with $\Pr_t=0.9$. Both are integrated over one pitch through the
  top of the near-wall layer, $z=z_x(x)$ of 4.2 ($z_x/D$ 0.59 at $x/D=4$ to 0.94 at 20), the same surface for
  both rungs, $x/D$ 3--20; vertical components only, the streamwise flux crossing the inclined surface is left
  out and stated. 193M uncertainty: jackknife over ten batches.
\item Findings: in the reference the turbulent flux exceeds the mean-flow flux from $x/D\approx5$ (4.5--4.75);
  in the RANS it never does. RANS over reference: turbulent flux a fraction (about half); mean-flow flux
  several times. Ratios to be confirmed once the flux surface is recomputed on the final $z_x$; the paper
  quotes no more than one ratio each. The RANS replaces the missing turbulent
  transport by a stronger mean upwash, through its stronger CRVP. The RANS centreline upwash velocity at the
  core height is 1.14, 1.66 and 1.20 times the reference at $x/D$ 4, 10, 20, and its inward velocity under the
  pair ($z/D=0.1$) 1.75--1.93 times.
\item Jet trajectory and spread: height of the coolant centroid and of the maximum of the temperature deficit on
  the centre plane, and the lateral half-width of the coolant, against $x/D$. The centroid height agrees within
  12\,\% (RANS 0.47, 0.69, 0.89\,$D$ against 0.47, 0.64, 0.79\,$D$ at $x/D$ 4, 10, 20). The RANS half-width is
  0.91, 0.80, 0.75 of the reference (0.87, 0.81, 0.91\,$D$ against 0.96, 1.02, 1.21\,$D$): the reference plume
  widens, the RANS plume does not. On the centre plane the coolest fluid of the reference is at the wall at every
  station; in the RANS it sits at $z/D$ 0.35--0.6 from $x/D=4$.
\item Consequence for $\eta$, in one paragraph, with the ratios RANS/193M at $x/D$ 5, 10, 20: plume half-width
  0.89, 0.80, 0.75; $\bar\eta$ 0.87, 0.70, 0.67; centreline $\eta$ 1.05, 1.18, 1.29. A stronger CRVP with a
  stronger mean upwash and inflow under the pair, and a weaker turbulent exchange, keeps the coolant on the
  centreline and lifts its core there, and sweeps hot gas under the flanks; this is the narrow, hot
  RANS footprint of (iv). The share of the $\eta$ difference that belongs to the mean flow against the share that
  belongs to the closure is taken up in 4.2 and tested in 4.3.
\item Literature: CRVP and anti-kidney vortices of shaped holes (Haven et al., 1997); mean secondary flow behind
  the 7-7-7 hole (Schroeder and Thole, 2016; Haydt and Lynch, 2019); over-prediction of the CRVP by
  eddy-viscosity models (Hoda and Acharya, 2000).
\end{itemize}

\textbf{Figure F8: what the instantaneous jet looks like.} Instantaneous centre-plane temperature of the 193M 3DNS
over the whole plate, and a near-hole view. The mean cross-planes are in F4.
\begin{itemize}
\item WHY: shows the instantaneous features that only the scale-resolving rungs hold (Kelvin--Helmholtz roll-up of
  the windward shear layer, the large structures at the edge of the film), against the mean picture of F4.
\item HOW: one record of the 193M ($tU_\infty/D=249.09$), centre plane $y=0$, non-dimensional temperature in
  \texttt{jet}, with the hole and the solid drawn. (a) $x/D$ $-3$ to 20, $z/D$ $-1$ to 3, equal aspect, with the
  stations of the F4 cross-planes marked. (b) $x/D$ $-1$ to 6, with the roll-up labelled. Format PNG.
\item WHAT: the shear layer rolls up from the windward lip and breaks down by $x/D\approx3$; the film thickens to
  $z/D$ 1--1.5 by $x/D$ 10--20 with large structures along its outer edge.
\end{itemize}
\planfig{fig_flow_features.png}{0.9}{F8 (\texttt{figures/code/fig\_flow\_features.py}): instantaneous 193M
centre-plane temperature, (a) $x/D$ $-3$ to 20, (b) near-hole view with the shear-layer roll-up.}

\textbf{Figure F8b: how the vortices move and how strong they are.} CRVP trajectory and circulation, and the
circulation of the WV, against $x/D$; RANS 50M against 193M 3DNS. Observation only.
\begin{itemize}
\item WHY: turns the pictures of F4 into numbers along $x/D$: the differences in the vortex system between the
  RANS and the reference, before any effect is drawn from them.
\item HOW: two panels against $x/D$ (3 to 20), RANS 50M dashed and 193M b010 solid, from the fine cross-planes at
  ten stations. (a) Core positions $|y_c|/D$ and $z_c/D$ of the CRVP (vorticity-weighted centroid inside half
  the peak; band between the two sides of the 193M). (b) Circulation $\Gamma/(U_\infty D)$ of a CRVP vortex and
  of the WV (area integral of $\omega_x$ inside the contour at 0.2 of the peak); open markers where a pick is
  below three times the plane noise or fails the swirl check. Format PDF.
\item WHAT: the RANS CRVP is lower by about 0.1\,$D$ and 1.2--1.8 times stronger; the reference keeps a WV to
  $x/D=15$, the RANS loses it by $x/D=5$.
\end{itemize}
\planfig{fig_mean_features.pdf}{0.9}{F8b (\texttt{figures/code/fig\_mean\_features.py}): (a) CRVP core position,
(b) circulation of the CRVP and of the wall vortex; RANS 50M and 193M b010.}

\textbf{Figure F8c: what these differences do to the coolant.} Coolant flux off the wall, and the chain to $\eta$;
RANS 50M against 193M 3DNS.
\begin{itemize}
\item WHY: the effect of the vortex differences of F8b: how the coolant leaves the wall layer in each rung, and
  how that appears in the width of the film and in $\eta$.
\item HOW: two panels against $x/D$ (4 to 20). (a) Mean-flow (upwash) and turbulent coolant flux through the top
  of the near-wall layer $z=z_x(x)$ of 4.2, vertical components, integrated over one pitch, both rungs (193M
  solid, RANS dashed), logarithmic axis, $\pm2$ standard errors (ten-batch jackknife) on the 193M. (b) Ratios
  RANS/193M of the lateral half-width of the coolant plume, of $\bar\eta$ and of the centreline $\eta$, with a
  line at 1; $\eta$ from the data of F11. Format PDF.
\item WHAT: in the reference the turbulent flux exceeds the mean-flow flux from $x/D\approx5$; in the RANS it
  never does. The RANS turbulent flux is a fraction of the reference (about half), in line with the near-wall
  eddy-viscosity deficit of F14; its mean-flow flux is several times the reference (ratios to be confirmed on
  the final $z_x$). By $x/D=20$ the RANS plume is 0.75 of the reference width, $\bar\eta$ 0.67 and the
  centreline $\eta$ 1.29.
\end{itemize}
\planfig{fig_flux_eta.pdf}{0.9}{F8c (\texttt{figures/code/fig\_flux\_eta.py}): (a) mean-flow and turbulent coolant
flux through the top of the near-wall layer $z_x(x)$, vertical part, (b) RANS/193M ratios of plume half-width,
$\bar\eta$ and centreline $\eta$.}

\textbf{(iii) The mainstream layer between the jets.}
\begin{itemize}
\item The approach layer is laminar under decaying FST: at $x/D=-2.3$ the 3DNS give $H$ 2.58--2.62 (Blasius
  2.59), $\mathrm{Re}_\theta$ 175--182; the RANS gives $H$ 2.88 and $\mathrm{Re}_\theta$ 165. FST streaks raise
  $u_{rms}$ in the layer to 0.07--0.10\,$U_e$ in the 3DNS (Jacobs and Durbin, 2001; Brandt et al., 2004).
\item Intermittency $\gamma$, one definition for every value from the 3DNS and LES in the paper: the detector is
  the high-pass spanwise stress $d=(v-\langle v\rangle_T)^2/U_\infty^2$, with $\langle v\rangle_T$ a running mean
  over $2\,D/U_\infty$, smoothed over $1\,D/U_\infty$; laminar and turbulent are separated by Otsu's method on
  $\log_{10}d$ (after Dellacasagrande et al., 2023, with the running mean in time in place of their spanwise
  mean, since this flow is not homogeneous in span); $\gamma$ is the fraction of time above the threshold. The
  threshold is taken from the 193M over the band where the two states coexist ($|y/D|$ 2--3, $x/D$ 0--20) and
  applied to both runs. Checks: $\gamma$ is $10^{-4}$ in the approach layer and 0.993 in the film core of the
  193M (LES 0.977, minimum 0.88). The RANS $\gamma$ is the transported variable of the model, a different
  quantity.
\item Behind the hole ($|y/D|\le1.5$) the jet trips the layer: the 3DNS layer is turbulent from $x/D\approx2$.
  The RANS layer behind the jet is turbulent too.
\item At mid-pitch the rungs part. In the 193M and the LES a laminar wedge lies between the turbulent wedges of
  neighbouring jets and closes over $x/D$ 10--15; turbulent patches reach the mid-pitch line from $x/D$ 5--6 as
  the wedge edges meander across it and as patches are shed ahead of the tip (evidence: F10). Mid-pitch $\gamma$
  of the 193M is 0.002 at $x/D=5$, 0.18 at 10, 0.74 at 15 and 0.94 at 20 (LES 0.17, 0.77, 0.94), first above 0.8
  at $x/D=15.7$; at $x/D=10$ it ranges over 0.10--0.29 when the threshold is doubled or halved, at $x/D=20$ over
  0.87--0.97. $H$ falls to 1.71. The RANS mid-pitch layer stays laminar to $x/D=20$ (model $\gamma$ 0.02,
  $\mu_t/\mu\le0.6$, $H$ 2.2--2.6, $C_f$ $1.8$--$2.2\times10^{-3}$ against $2.4$--$5.6\times10^{-3}$). No spot
  celerities are reported: the patches on the mid-pitch line are mostly wedge edges, and too few isolated spots
  remain for a statistic.
\item Consequence: near-wall $k$ of the RANS between the jets is 0.02--0.14 of the 3DNS value, so there is no
  turbulent mixing to carry coolant to the mid-pitch.
\item Literature: holes trip a laminar layer (Goldstein and Yoshida, 1982); spot growth and merger in bypass
  transition (Zaki, 2013; Ovchinnikov et al., 2008); $\gamma$--$\mathrm{Re}_\theta$ has no film-cooling trip
  (Langtry and Menter, 2009).
\item Instantaneous features seen in the 193M: the Kelvin--Helmholtz roll-up of the jet shear layer (F8 g) and
  the turbulent patches between the jets (F10). Steady RANS carries the mean only, so it has neither.
\item Share of the mixing carried by the coherent features: the periodic (coherent) share of
  $\mu_{\mathrm{eff}}$ is 0.05--0.08 in the hole, jet and film at a filter width $\sigma$ of 0.036 (method of Lav
  et al., 2019; 0.03--0.06 at their width), on the 193M. The mixing is mostly broadband turbulence, as Fawcett
  (2011) predicts at VR near 1.
\end{itemize}

\textbf{Figure F7: how $H$ and $C_f$ develop along the plate.} $H$ and $C_f$ against $x/D$ on the centreline
and at mid-pitch, six rungs.
\begin{itemize}
\item WHY: puts the state of the layer into the parameters the reader knows, and shows where each rung departs.
  First figure of the mainstream chain.
\item HOW: $H$ (a, b) and $C_f$ (c, d) against $x/D$, centreline and mid-pitch, for the six rungs. The Blasius
  and turbulent correlations are labelled by text and arrows in the plot, not in the legend, which lists the
  cases only.
\item WHAT: the 20M LES shows a spike at $x/D$ 19.8--20, which is its outlet. The laminar RANS mid-pitch values
  of (iii) ($H$ 2.2--2.6, $C_f$ $1.8$--$2.2\times10^{-3}$) lie against the 3DNS values of $H$ 1.71 and $C_f$
  $2.4$--$5.6\times10^{-3}$.
\end{itemize}
\planfig{fig_bl_params.pdf}{0.95}{F7 (\texttt{fig\_bl\_params.py}): $H$ (a, b) and $C_f$ (c, d), centreline and
mid-pitch, six rungs (193M b010, others b006); Blasius and turbulent correlations labelled in the plot.}

\textbf{Figure F9: where the wall layer is turbulent.} Plan view of the plate, near-wall $k/U_\infty^2$ and
$\gamma$; RANS 50M, 20M LES, 193M 3DNS.
\begin{itemize}
\item WHY: shows the turbulent wedge behind each jet and the spread transition between the jets, against the
  laminar RANS mid-pitch.
\item HOW: plan view of the plate ($x/D$ $-3$ to 20, one pitch). Top row: $k/U_\infty^2$ at $z/D=0.1$ (RANS 50M
  cell layer, 20M LES b006, 193M b010). Bottom row: $\gamma$ (RANS: model variable; 20M LES and 193M: detector
  of (iii) on their own $z/D=0.1$ samplers, one threshold). Format PNG, \texttt{inferno}.
\item WHAT: mean $\gamma$ over the plate is 0.17 for the RANS (turbulent only behind the jet, about 0.02 between
  the jets), 0.73 for the 20M LES and 0.72 for the 193M.
\end{itemize}
\planfig{fig_wall_state.png}{0.95}{F9 (\texttt{fig\_wall\_state.py}): $k/U_\infty^2$ at $z/D=0.1$ (top) and $\gamma$
(bottom); RANS 50M, 20M LES b006, 193M b010; $x/D$ $-3$ to 20, one pitch.}

\textbf{Figure F10: how the layer between the jets becomes turbulent.} The turbulence detector of (iii) on the
$z/D=0.1$ plane, 193M 3DNS (b010) and 20M LES: snapshots, a space--time diagram at mid-pitch, and $\gamma(x)$.
\begin{itemize}
\item WHY: the evidence for the spread transition: a laminar wedge between the turbulent wedges of neighbouring
  jets, which closes downstream, with turbulent patches crossing the mid-pitch line before it closes.
\item HOW: left column 193M, right column 20M LES. Top: four plan-view snapshots of $\log_{10}d$ on the
  $z/D=0.1$ plane (two pitches, mid-pitch line at $y/D=3$), with the threshold contour, the hole exits, and one
  tracked patch marked. Middle: $(x,t)$ diagram of $\log_{10}d$ along the mid-pitch line over 40\,$D/U_\infty$,
  time measured from the start of each window, snapshot times marked. Bottom: $\gamma(x)$ at mid-pitch for the
  193M and the LES, each with a band giving $\gamma$ when the threshold is halved (upper edge) and doubled (lower
  edge), and the model $\gamma$ of the RANS 50M, labelled as such. One colour scale, \texttt{inferno}. Format PNG.
\item WHAT: the wedge closes over $x/D$ 10--15 in both runs; the 193M and the LES agree in $\gamma(x)$ within 0.03
  (0.18, 0.74, 0.94 against 0.17, 0.77, 0.94 at $x/D$ 10, 15, 20); the RANS model $\gamma$ stays at 0.02. No
  celerities.
\end{itemize}
\planfig{fig_spots.png}{0.95}{F10 (\texttt{fig\_spots.py}): turbulence detector $d$ (high-pass spanwise stress) on
the $z/D=0.1$ plane, 193M b010 (left) and 20M LES (right): snapshots with the threshold contour, $(x,t)$ diagram
at mid-pitch, and $\gamma(x)$ with threshold bands and the RANS model $\gamma$.}

\textbf{(iv) The outcome: effectiveness.}
\begin{itemize}
\item $\bar\eta$ over $x/D$ 2.5--7 / 7--15 / 15--20: RANS 0.270--0.272 / 0.168--0.171 / 0.138 (mesh-independent);
  20M LES 0.301 / 0.231 / 0.202; 50M 0.305 / 0.240 / 0.209; 100M 0.307 / 0.239 / 0.211; 193M (b010) 0.304 /
  0.237 / 0.208; experiment at Tu 0.5\,\% 0.293 / 0.243 / 0.212. The scale-resolving rungs agree within 0.008; the RANS is
  29--33\,\% low beyond $x/D=7$.
\item Centreline against mid-pitch: the RANS centreline is too high ($\eta$ 0.497 against 0.387 for the 193M at
  $x/D=20$), its mid-pitch too low (0.005--0.014 against the 193M 0.025 at $x/D=20$; experiment
  0.033--0.042).
\item Three things feed this: the hole sets the jet (BR); the mean flow of (ii) sets where the CRVP holds and
  lifts the coolant; the layer of (iii) sets the lateral mixing (a laminar RANS mid-pitch carries no coolant
  across the pitch). The RANS film is too narrow, not too
  weak.
\item The 3DNS mid-pitch $\eta$ is below the experiment (193M 0.004, 0.007, 0.015, 0.025 at $x/D$ 5, 10, 15,
  20 against 0.037, 0.033, 0.033, 0.035), although $\bar\eta$ agrees. Two explanations, both supported, neither
  alone sufficient:
  \begin{itemize}
  \item The approach layer differs. The experiment trips its layer at $x/D=-35$: turbulent, $H$ 1.45,
    $\mathrm{Re}_\theta$ 670 at $x/D=-4.7$ (Schroeder and Thole, 2014). The 3DNS layer is laminar, $H\approx2.6$,
    $\mathrm{Re}_\theta\approx180$. A turbulent layer mixes coolant to the mid-pitch from the hole on (the
    measured mid-pitch $\eta$ is flat at 0.033--0.037 from $x/D=2$).
  \item Transition takes distance. Tested on the 193M: $\eta_{\mathrm{mid}}$ is flat at 0.004--0.007 to
    $x/D=9$ and rises with $\gamma_{\mathrm{mid}}$; the gap narrows from $-0.033$ at $x/D=5$ to $-0.010$ at
    $x/D=20$. The mid-pitch $\gamma$ exceeds 0.8 from $x/D=15.7$; at $x/D=20$, inside that stretch, the 3DNS is
    still 0.010 below the experiment (0.025 against 0.035), so a fully turbulent mid-pitch does not close the
    gap: the coolant that a turbulent approach layer carries across from the hole on is missing upstream.
    Nothing is reported beyond $x/D=20$.
  \end{itemize}
  So the gap belongs to the laminar approach layer of the configuration (transition delay plus the missing
  upstream lateral mixing), not to a model or resolution error: all scale-resolving rungs agree (mid-pitch
  $\eta$ at $x/D=20$ 0.025--0.029). F11 (c) carries $\gamma_{\mathrm{mid}}(x)$ on a secondary axis.
\end{itemize}

\textbf{Figure F11: how effectiveness varies with $x/D$.} (a) $\bar\eta$, (b) centreline $\eta$, (c) mid-pitch
$\eta$; six rungs and the 7-7-7 data at DR 1.5, BR 1.5.
\begin{itemize}
\item WHY: the result the paper explains: the scale-resolving rungs on the data, the RANS low in $\bar\eta$,
  high on the centreline and near zero at mid-pitch.
\item HOW: the six rungs and the experiment (DR 1.5, BR 1.5, Tu 0.5\,\%) are plotted against $x/D$ (E5). Panel
  (c) is taken at $|y/D|=2.91$ and carries the 193M $\gamma_{\mathrm{mid}}$ on a secondary axis (dashed grey),
  with a band for the detector threshold halved and doubled.
\item WHAT: the numbers of (iv). The scale-resolving rungs agree within 0.008 in $\bar\eta$; the RANS is 29--33\,\%
  low beyond $x/D=7$, with a centreline $\eta$ of 0.497 against 0.387 at $x/D=20$ and a mid-pitch $\eta$ of
  0.005--0.014 against 0.025.
\end{itemize}
\planfig{fig_eta_x.pdf}{0.55}{F11 (\texttt{fig\_eta\_x.py}): (a) $\bar\eta$, (b) centreline $\eta$, (c) mid-pitch
$\eta$ ($|y/D|=2.91$) with the 193M $\gamma_{\mathrm{mid}}$ (dashed grey, secondary axis); 193M b010, other rungs
b006, experiment (DR 1.5, BR 1.5, Tu 0.5\,\%).}

\textbf{Figure F12: where the coolant lands on the wall.} $\eta$ on the plate, $x/D$ 0--20, one pitch; RANS 50M,
20M LES, 193M 3DNS.
\begin{itemize}
\item WHY: shows the narrow RANS footprint against the wide resolved one, and links F9 (wall state) to F11.
\item HOW: $\eta$ (E5) is drawn on the plate for (a) RANS 50M, (b) 20M LES and (c) 193M 3DNS. Format PNG,
  \texttt{jet}.
\item WHAT: the RANS footprint is narrow and the resolved footprints are wide, which is the centreline and
  mid-pitch difference of F11 seen on the wall. The exit footprint is blank, since no plate faces lie there.
\end{itemize}
\planfig{fig_eta_wall.png}{0.75}{F12 (\texttt{fig\_eta\_wall.py}): $\eta$ on the plate, (a) RANS 50M, (b) 20M LES, (c)
193M 3DNS (b010; others b006).}

\textbf{(v) From the flow to the closure.}
\begin{itemize}
\item The coolant reaches the wall through the turbulent stress and heat flux of the mean equations. To attribute
  the $\eta$ differences to the flow, the two inputs of the RANS heat flux are measured in the reference: the
  eddy viscosity, as an effective eddy viscosity $\mu_{\mathrm{eff}}$ (E9), and the turbulent Prandtl number in
  its classic form (E10). The RANS heat flux is $c_p\mu_t/\Pr_t\,\nabla\tilde T$ (E2), so it is wrong if either
  input is wrong.
\item For the RANS they equal $\mu_t$ and 0.9 by construction (E1, E2); for the 3DNS they are measured from the
  stresses and fluxes. 4.2 breaks the $\mu_t$ difference down and compares the $\Pr_t$.
\item No effective diffusivity $\alpha_{\mathrm{eff}}$ and no tensor form of $\Pr_t$ appear anywhere in the
  paper; one linear story.
\end{itemize}
\begin{equation}
\mu_{\mathrm{eff}}=-\frac{a\!:\!S^{*}}{2|S^{*}|^2},\qquad a^{\perp}=a+2\mu_{\mathrm{eff}}S^{*},\qquad a^{\perp}\!:\!S^{*}=0
\tag{E9}
\end{equation}
Purpose: the eddy viscosity that best fits the measured anisotropy $a$ to the mean strain $S^{*}$; $a^\perp$ is
the part no eddy viscosity can carry.
\begin{equation}
\Pr_t=\frac{\widetilde{u''w''}\big/\partial\tilde u/\partial z}{\widetilde{w''T''}\big/\partial\tilde T/\partial z},\qquad
\Pr_{t,R}=\frac{\sum_R\widetilde{u''w''}\,\partial_z\tilde u\big/\sum_R(\partial_z\tilde u)^2}
{\sum_R\widetilde{w''T''}\,\partial_z\tilde T\big/\sum_R(\partial_z\tilde T)^2}
\tag{E10}
\end{equation}
Purpose: the $\Pr_t$ of the paper, the quantity E2 sets to 0.9, in the form of Kohli and Bogard (2005); the
regional value $\Pr_{t,R}$ is the gradient-weighted least-squares fit, so low-gradient points carry little weight
and no point thresholds are needed; quoted for $z/D\ge0.3$, since the gradient weighting lets the viscous wall
layer dominate otherwise.

\textbf{Figure F13: what the turbulent Prandtl number of the reference is.} Classic $\Pr_t$ (E10) of the 193M
3DNS, averaged over the span: map over $x/D$ and $z/D$.
\begin{itemize}
\item WHY: first look at the closure: E2 assumes 0.9, and the figure shows what the reference has in the film.
  Opens 4.2.
\item HOW: span-summed least-squares $\Pr_t$ (E10) of the 193M b010 over one pitch, drawn at the resolution of
  the mesh, \texttt{inferno}, 0 to 2, with 0.9 marked on the colour bar. Grey where $\Pr_t$ is not defined: the
  root-mean-square gradient of velocity or of temperature is below $0.05\,U_\infty/D$ or $0.05\,\Delta T_0/D$, or
  the heat diffusivity is within two standard errors of zero (ten batches). The dashed line is the boundary of the
  defined region.
\item WHAT: in the film the reference mixes heat faster than momentum where E2 assumes 0.9: 0.5--0.9 in the body
  of the film, below 0.3 along the jet core at $z/D$ 0.4--0.7. Span-averaged values with their standard errors,
  near film: $0.39\pm0.05$, $0.18\pm0.01$, $0.12\pm0.01$ at $z/D$ 0.2, 0.3, 0.5; far film: $0.89\pm0.10$,
  $0.83\pm0.09$, $0.46\pm0.04$. The near film is resolved down to $z^+=2.5$. In the far film the layer below
  $z^+\approx13$ is not resolved: the wall is adiabatic, the film is mixed, and both the heat flux and the
  temperature gradient are within the sampling noise, so their ratio is not determined. No counter-gradient layer
  is claimed.
\end{itemize}
\planfig{fig_prt_classic_span.png}{0.9}{F13 (\texttt{fig\_prt\_classic\_span.py}): span-averaged classic $\Pr_t$
(E10) of the 193M b010; grey where it is not defined; dashed line: boundary of the defined region; 0.9 marked on
the colour bar.}

%-----------------------------------------------------------------------------------------------------
\subsection*{4.2\quad Breakdown of Model Differences}
\textbf{Purpose.} One linear story. The RANS heat flux is $c_p\mu_t/\Pr_t\,\nabla\tilde T$
(E2), so it has two inputs. (1) $\mu_t$ against the point-wise optimal eddy viscosity of the reference (E9), seen
as span-averaged maps over $x/D$ 3--20 (the jet, $x/D<3$, is left out). Points are kept only where the
reference strain satisfies $|S^*|D/U_\infty>0.05$, point by point in three dimensions before the span average.
``Near wall'' is one layer, used throughout the paper: from the wall to $z_x(x)$, the height at which, going up
from the wall, the $\mu$ ratio RANS/193M first reaches one ($z_x/D$ 0.59 at $x/D=4$ to 0.94 at 20). Below
$z_x$ the RANS $\mu_t$ is about 0.55--0.66 of the reference, mostly through a short length scale, with $k^{1/2}$
also low and the assumed alignment giving back part. Above $z_x$ the two are close, through compensating errors
in $k^{1/2}$, length and alignment. (2) The heat flux inherits that, and the
$\Pr_t$ setting of 0.9 is higher than the reference value: the reference classic $\Pr_t$, fitted over the volume of each region, is 0.30 in the jet, 0.23 in the near film and 0.45 in the far film, at or below the free-shear value
(Kays, 1994; Schlichting and Gersten, 2017). (3) Spanwise: on the film flanks the reference heat diffusivity is larger
laterally than wall-normally, while its momentum diffusivity is about isotropic; the RANS lateral heat
diffusivity there is about half the reference (F19). Each point is compared with the literature; what is new is the
split and the assigned causes. Ends on the remedies, which lead into 4.3. About 3 pages.

\textbf{(i) The momentum factors and the log budget.}
\begin{itemize}
\item $\mu_{\mathrm{eff}}$ is the product of how much turbulence there is ($k^{1/2}$), the size of the motions
  that mix ($L_{\mathrm{eff}}$) and how well the stress lines up with the strain ($\cos\theta$) (E11).
\item The log of a ratio of products is a sum, so the difference between a rung and the reference splits
  exactly into the factors at each point (E12).
\item Length and alignment are not independent: $L_{\mathrm{eff}}=L_{\mathrm{al}}/\cos\theta$ with the aligned length
  $L_{\mathrm{al}}=\mu_{\mathrm{eff}}/(\bar\rho k^{1/2})$, so the two terms share $\ln(1/\cos\theta^{193})$ with opposite
  signs and only their sum $\ln(L_{\mathrm{al}}/L_{\mathrm{al}}^{193})$ is determined. The paper reports $L_{\mathrm{al}}$
  as the term, and $L_{\mathrm{eff}}$ and $\cos\theta$ separately only with their attribution ranges.
\item The estimator is the point-wise optimal eddy viscosity of the literature (E9; Wu et al., 2018; McConkey et
  al., 2022). It is shown as conditional span means over $(x,z)$ for the RANS and the reference (F14).
\item Mask: $\mu_{\mathrm{eff}}$ divides by $|S^*|^2$, so it has no meaning where there is no strain. Points are
  kept only where the 193M strain satisfies $|S^*|D/U_\infty>0.05$, applied point by point in three dimensions
  before the span average (a masked comparison as in Stratton and Shih, 2019). A map cell is drawn only where at
  least $f_{\min}=0.2$ of its span points pass; elsewhere it is grey. Range $x/D$ 3--20, $z/D$ 0--2; the jet
  ($x/D<4$) is left out of the comparison.
\item Near wall: the layer from the wall (wall nodes excluded: $k=0$ on both rungs) to $z_x(x)$, the height at
  which, going up from the wall, the ratio RANS/193M of the conditional span-mean eddy viscosity first reaches
  one. $z_x/D$ rises from 0.59 at $x/D=4$ to 0.94 at $x/D=20$. One definition for F14, F14b and the coolant flux
  of 4.1 (F8c); the same surface for every rung.
\item Value of the layer (F14b): in slabs of $1\,D$ along $x/D$, the log ratio of the volume means of $\mu$ over the
  masked points of the layer. The four terms (density, $k^{1/2}$, $L_{\mathrm{eff}}$, $\cos\theta$) are formed
  from the point values and add exactly to the total, without an order to choose.
\item The energy factor ties to the familiar $k$ budget (E14): production $2|S^*|b_\parallel$ against
  dissipation; the length ties to the dissipation length: on the 193M centreline $L_{\mathrm{eff}}/(C_\mu
  k^{3/2}/\varepsilon)$ is 1.02--1.07 (F18), and for SST $\mu_t/(\bar\rho k^{1/2})=C_\mu k^{3/2}/\varepsilon$ exactly.
\item The RANS has $\cos\theta=1$ and $\Pr_t=0.9$ by construction; the split shows what that costs. The alignment
  test of the Boussinesq hypothesis follows Schmitt (2007).
\end{itemize}
\begin{equation}
\mu_{\mathrm{eff}}=\bar\rho\,k^{1/2}\,L_{\mathrm{eff}}\cos\theta,\qquad L_{\mathrm{eff}}=\frac{k^{1/2}|b|}{|S^*|}
\tag{E11}
\end{equation}
\begin{equation}
\ln\frac{\mu_{\mathrm{eff}}}{\mu_{\mathrm{eff}}^{193}}=\ln\frac{\bar\rho}{\bar\rho^{193}}
+\underbrace{\tfrac12\ln\frac{k}{k^{193}}}_{\text{energy}}
+\underbrace{\ln\frac{L_{\mathrm{eff}}}{L_{\mathrm{eff}}^{193}}}_{\text{length}}
+\underbrace{\ln\frac{\cos\theta}{\cos\theta^{193}}}_{\text{alignment}}
\tag{E12}
\end{equation}
\begin{equation}
\mu_R=\frac{1}{V_R}\int_R\mu\,\mathrm dV\ \ \text{(masked points of the layer slab)},\qquad
\ln\frac{\mu_R}{\mu_R^{193}}=\phi_\rho+\phi_k+\phi_{L}+\phi_{\cos\theta}
\tag{E13}
\end{equation}
\begin{equation}
\mathrm D\ln k=2|S^*|\,b_\parallel-\tfrac23\tilde S_{kk}-\frac{\varepsilon}{\bar\rho k}+\mathcal T_k,\qquad b_\parallel=|b|\cos\theta
\tag{E14}
\end{equation}
Purpose of E11--E14: E11 names the momentum factors; E12 makes them additive per point; E13 defines the
layer value and its exact split ($\phi$ terms), which is what F14b shows; E14 links the energy factor
to production and dissipation.

\textbf{(ii) Region by region.}
\begin{itemize}
\item Order: the near-wall layer below $z_x$, then the strained film above it, over $x/D$ 3--20. The jet
  ($x/D<4$) is left out. Hole and approach layer one sentence each: the hole values are wall-layer values (over
  99.9\,\% of the metering strain weight within 0.1\,$D$ of the wall).
\item Rungs in one sentence: the ratio to the reference is the same for the RANS 5M and 50M, so the RANS
  difference is the model and not the mesh; the 20M LES and 50M 3DNS lie about 10\,\% above the reference.
\item \textbf{Step 1, $\mu_t$ (F14, F14b, F16, F17, F18).} RANS 50M against the 193M b010.
\item Finding 1, below $z_x$ (F14 d, F14b a): the RANS $\mu_t$ is about 0.55--0.66 of the reference (``about a
  half to two thirds'' in the abstract and conclusions); $\ln$ total $-0.42$, $-0.55$, $-0.59$ at $x/D$ 5, 10,
  20.
\item Finding 2, the split (F14b): the length term is the largest deficit, RANS $L$ about 0.45--0.5 of the
  reference; $k^{1/2}$ about 0.6; density negligible. Terms: $L_{\mathrm{eff}}$ $-0.67$ to $-0.74$, $k^{1/2}$
  $-0.22$ to $-0.30$, $\cos\theta$ $+0.35$ to $+0.46$, density $+0.01$ to $+0.05$. The perfect alignment the RANS assumes gives back part:
  $1/\cos\theta$ of the reference 2.5 near the hole, falling to 1.7 at $x/D=20$. Without the alignment assumption
  the near-wall deficit would be larger still. Stated neutrally; terms of opposite sign are reported together.
\item Finding 3, above $z_x$ (F14): the $\mu$ ratio is close to one (1.0--1.12, mean 1.06), through
  compensating errors: $k^{1/2}$ 1.10--1.16, length 0.77--0.85, $1/\cos\theta$ 1.23--1.45. The paper writes ``close
  to one'' with at most one number. A compensation of this kind has not been found described for film cooling;
  the nearest works are those already cited (optimal eddy viscosity: Wu et al., 2018; McConkey et al., 2022;
  alignment: Schmitt, 2007; Stratton and Shih, 2019; Sandberg et al., 2018). No novelty claim beyond that.
\item Above the strained film: grey in F14, no comparison. The RANS carries a finite $\mu_t$ there; for heat it
  still acts through $\mu_t/\Pr_t$ wherever a temperature gradient exists, which is the subject of step 2 and of
  the coolant flux of 4.1.
\item Where the alignment term comes from: the film volume median $\cos\theta$ of the 193M is 0.62 (near) and 0.72
  (far); the centre plane is the best-aligned part (median 0.83); misalignment sits near the wall and on the film
  flanks (F17). The energy term is largest in the far film, where the RANS mid-pitch layer is laminar (4.1 (iii)).
\item \textbf{Step 2, $\Pr_t$ (F13, F15).} The RANS heat flux $c_p\mu_t/0.9\,\nabla\tilde T$ carries the $\mu_t$
  deficit of step 1, divided by a $\Pr_t$ that is higher than the reference value.
\item Origin of 0.9: the solver default (stated in the text, no citation), representative of the 0.85--0.9
  backed out of log-layer profiles of a flat-plate boundary layer (Kays, 1994); about 0.5 is the free-shear value
  (Schlichting and Gersten, 2017); Kohli and Bogard (2005) note most models assume $\Pr_t\approx1$.
\item The setting (193M b010, classic form, least squares over the volume of each region, full pitch,
  $z/D\ge0.3$, standard error over ten batches): jet $0.300\pm0.004$, near film $0.225\pm0.005$, far film
  ($x/D$ 10--20) $0.454\pm0.013$. Below 0.9 in every region, and at or below
  the free-shear value of about 0.5. The whole simulated volume is used, since the RANS applies its setting
  everywhere. On the centre plane alone the values are higher, 0.44, 0.63 and 0.80 (jet, near film, far film to
  $x/D=20$): the centre plane is where $\Pr_t$ is highest, and the film flanks lower the regional value.
\item Near the wall (span-averaged, F13): the near film is resolved to $z^+=2.5$, with $\Pr_t$ $0.39\pm0.05$ at
  $z/D=0.2$ and $0.81\pm0.21$ at 0.1. The far film carries a layer near 0.9 at $z^+$ 13--40 ($0.83$--$0.98$,
  standard error 0.09--0.18) and is not resolved below $z^+\approx13$. No converged $\Pr_t\approx2$ (high
  confidence).
\item Consequence: lowering $\Pr_t$ to the reference value corrects the heat flux for a given $\mu_t$, but not the
  $\mu_t$ deficit of step 1, nor the lateral heat anisotropy of step 3.
\item \textbf{Step 3, spanwise (F19).} On the film flanks the reference heat diffusivity is anisotropic (lateral
  larger than wall-normal); the momentum diffusivity is about isotropic. Numbers (193M b010, cross-planes, regional least squares, $z/D\ge0.3$; 95\,\% in brackets):
  $\alpha_y/\alpha_z$ 3.2 (2.3--4.4) on the film flanks ($|y|/D\ge1.5$) at $x/D=10$ and 1.7 (1.4--1.9) at
  $x/D=20$, inside or near the 1.5--4 of Kaszeta and Simon (2000); about 1 in the plume centre. The momentum ratio
  $\nu_y/\nu_z$ is about 1 on the flanks (0.7--1.1); the heat anisotropy does not follow from the momentum.
\item Lateral heat diffusivity at the film edges ($|y|/D\ge1$), RANS 50M against 193M: 0.54 times at $x/D=10$,
  0.63 times at $x/D=20$. The lateral $\Pr_t$ of the reference there is 0.44 and 0.56, against 0.9: the setting
  is the largest term of that deficit; the RANS $\mu_t$ against the reference $\mu_z$ nearly cancels the
  anisotropy term. The RANS lateral flux integral still matches the 193M, because its lateral gradients are
  steeper (squared ratio 1.9): the narrow, steep RANS film is the result of the weak lateral diffusivity.
  Confidence medium (edge definition, finite-difference RANS gradient); totals b006 and b010 within 0.05.
\item Known in the literature as the RANS culprit (Acharya et al., 2001; Harrison and Bogard, 2008a); new here:
  the measured lateral diffusivity deficit assigned to the $\Pr_t$ setting and the anisotropy, for heat. Not shown:
  a momentum anisotropy at the film edges.
\item RANS sweep link (reference 193M b010): global $\Pr_t$ 0.5, 0.4,
  0.3 move the centreline $\eta$ at $x/D=20$ from 0.497 to 0.361 (193M 0.387) but $\bar\eta$ ($x/D$ 15--20)
  saturates at 0.160--0.164 (193M 0.208); QCR gives 0.185, QCR with $\Pr_t$ 0.4 0.198. A $\Pr_t$ setting fixes the
  centreline, not the lateral spread; QCR alone raises the centreline (0.524). Recovered share of the $x/D$ 7--20
  $\bar\eta$ deficit against the 193M: $\Pr_t$ 0.5 / 0.4 / 0.3 alone 36 / 44 / 48\,\%; QCR alone 72\,\%; QCR with
  $\Pr_t$ 0.4 97\,\%; local $\Pr_t$ (0.4 near the hole, 0.9 beyond $x/D\approx10$) 34\,\%, with QCR 96\,\%. The
  stress form is worth more than the $\Pr_t$ setting. Mid-pitch $\eta$ at $x/D=20$ stays 0.006--0.012 in every
  $\Pr_t$ and QCR run (193M 0.025). Leads to 4.3, which adds two runs: P1, zonal
  $\Pr_t$ (0.4 for $x/D\ge3$, 0.9 in the jet and hole), and P2, heat-flux substitution in the film only; both
  50M SST. The laminar RANS mid-pitch is tested by RC-S9: $\gamma=1$ for
  $x>0$, $z>0$, approach layer and hole unchanged; 5M: $\eta_{\mathrm{plate}}$ $+1.3$\,\%, 50M smoke test $+1.26$\,\%.
\item Kohli and Bogard (2005): $\Pr_t\approx2$ at low freestream turbulence (0.6\,\%, a factor 4 variation) and
  $\approx1$, nearly uniform, at high (20\,\%). The 193M freestream $Tu$ is 3.3--3.6\,\% ($x/D$ $-5$ and 0,
  $z/D$ 2--3), between the two. The present near-wall 1.0--1.3 at $z/D=0.1$ falling to 0.5--0.6 sits nearer their
  high-turbulence value; their geometry (round holes, DR 1.05) also differs.
\item What the mixing motions are: the periodic (coherent) share of $\mu_{\mathrm{eff}}$ is 0.05--0.08 in the
  hole, jet and film at a filter width $\sigma$ of 0.036 (method of Lav et al., 2019; 0.03--0.06 at their width);
  the mixing is mostly broadband turbulence, as Fawcett (2011) predicts at VR near 1, and the RANS gets the size
  of these motions wrong.
\item Alignment in the reference: mean strain builds the aligned anisotropy and mean rotation removes 55--114\,\%
  of it; a linear eddy viscosity keeps the strain drive only (budgets of $b_\parallel$, $b_\perp$).
\item Literature at each point: SST $\nu_t$ too low in the film (Stratton and Shih, 2019); non-uniform,
  anisotropic diffusivity (Kaszeta and Simon, 2000, ratio 1.5--4; Coletti et al., 2013); low and varying $\Pr_t$
  (Ling et al., 2016; Liu et al., 2008, 2011; Milani et al., 2019); Kohli and Bogard (2005) as above; lateral
  temperature spread under-predicted with $\Pr_t=0.9$ in a planar jet, by about 20\,\% (Rochhausen et al., 2015).
  For each: agrees or not, in one clause.
\end{itemize}

\textbf{Figure F14: where the RANS eddy viscosity differs from the reference, and in which factor.} Four
stacked maps of the RANS 50M over 193M ratio of the conditional span means.
\begin{itemize}
\item WHY: the RANS heat flux is proportional to $\mu_t$ (E2). The figure shows where the RANS $\mu_t$ is below
  the eddy viscosity of the reference and where it is not, and which factor makes the difference.
\item HOW: $x/D$ 3--20, $z/D$ 0--2; ratio RANS/193M of the conditional span means (points with 193M
  $|S^*|D/U_\infty>0.05$, 3D mask before the span average) of (a) $k^{1/2}$, (b) the length (RANS
  $\mu_t/(\bar\rho k^{1/2})$ over the 193M $L_{\mathrm{eff}}$), (c) $1/\cos\theta$ (RANS $\cos\theta=1$), (d) $\mu$
  ($\mu_t$ over $\mu_{\mathrm{eff}}$ of E9). One black line on every panel: $z_x(x)$. Grey: fewer than
  $f_{\min}=0.2$ of the span points pass the mask. Logarithmic colour scale 0.1 to 10, \texttt{bwr}, blue below
  one, red above. Format PNG.
\item WHAT: below $z_x$ blue in (d), with (b) the deepest blue and (c) red; above $z_x$ (d) near one while (a)
  is light red, (b) light blue and (c) red. $z_x/D$ rises from 0.59 at $x/D=4$ to 0.94 at 20.
\end{itemize}
\planfig{fig_meff_maps.png}{0.95}{F14 (\texttt{fig\_meff\_maps.py}): ratio RANS 50M / 193M b010 of the conditional
span means of (a) $k^{1/2}$, (b) length, (c) $1/\cos\theta$, (d) $\mu$; black line $z_x(x)$; grey where fewer than
20\,\% of the span points pass the strain mask.}

\textbf{Figure F14b: the factor split of the near-wall deficit along the plate.} Two panels against $x/D$ for
the near-wall layer.
\begin{itemize}
\item WHY: gives the size of each factor in the near-wall layer as a continuous curve, with an exact split.
\item HOW: slabs of $1\,D$, centres every $0.25\,D$, $x/D$ 3--20; layer from the wall to $z_x(x)$, masked
  points only. (a) $\ln(\mu_{\mathrm{RANS}}/\mu_{193})$ of the layer (E13) and its four terms: density,
  $k^{1/2}$, $L_{\mathrm{eff}}$, $\cos\theta$; the terms add exactly to the total. (b) Ratios RANS/193M of
  $k^{1/2}$, $L$ and $1/\cos\theta$. Format PDF, single column.
\item WHAT: total ratio about 0.55--0.66; length the largest deficit ($L$ about 0.45--0.5), $k^{1/2}$ about 0.6,
  density negligible; $1/\cos\theta$ 2.5 near the hole falling to 1.7 at $x/D=20$.
\end{itemize}
\planfig{fig_meff_split_film.pdf}{0.5}{F14b (\texttt{fig\_meff\_split\_film.py}): near-wall layer, (a) exact split
of $\ln(\mu_{\mathrm{RANS}}/\mu_{193})$ into density, $k^{1/2}$, $L_{\mathrm{eff}}$ and $\cos\theta$ terms, with the
total; (b) ratios RANS/193M of $k^{1/2}$, $L$ and $1/\cos\theta$.}

\textbf{Figure F15: how the $\Pr_t$ of the reference compares with the setting and the literature.} Classic
$\Pr_{t,R}$ (E10) of the 193M b010, fitted over the volume of the jet, near film and far film for $z/D\ge0.3$.
\begin{itemize}
\item WHY: step 2 in one view: the setting of 0.9 is higher than the reference value in every region.
\item HOW: one panel with the regional least-squares $\Pr_{t,R}$ (E10) of the 193M b010 over the volume of each
  region (filled); the centre-plane values beside them (open) show how the single plane differs from the
  region. Lines mark the RANS 0.9, the log-layer 0.85 (Kays, 1994), the free-shear 0.5 (Schlichting and Gersten,
  2017; Rochhausen et al., 2015) and Kohli and Bogard (2005), $\approx1$ and $\approx2$; each line is labelled in
  the plot.
\item WHAT: the volume values are $0.300\pm0.004$ (jet), $0.225\pm0.005$ (near film) and $0.454\pm0.013$ (far
  film), against 0.9; all are at or below the free-shear value. The centre-plane values are 0.44, 0.63 and 0.80.
  Bars of $\pm2$ standard errors over ten batches.
\end{itemize}
\planfig{fig_prt_regions.pdf}{0.6}{F15 (\texttt{fig\_prt\_regions.py}): classic regional $\Pr_{t,R}$ (E10), $z/D\ge0.3$,
193M b010: volume of each region (filled) and centre plane (open); jet, near film, far film. Lines: RANS 0.9,
log layer 0.85, free shear 0.5, Kohli and Bogard $\approx1$ and $\approx2$.}

\textbf{Figure F16: where the RANS has less turbulence and shorter mixing motions.} $y$--$z$ planes at $x/D$ 4,
10 and 20: $k/U_\infty^2$ (top) and $L_{\mathrm{eff}}/D$ (bottom), RANS 50M against 193M 3DNS, with $\eta$
contour lines.
\begin{itemize}
\item WHY: the flow field behind F14: where the RANS has less turbulence (mid-pitch, far film) and where its
  motions are too small (jet and film edges), next to where the coolant sits.
\item HOW: $k/U_\infty^2$ (a--f) and $L_{\mathrm{eff}}/D$ (g--l) (E11) on cross-planes at $x/D$ 4, 10, 20, for the
  RANS 50M and the 193M 3DNS (b010). White contours: $\eta$ 0.1--0.4. No masking: every point is drawn.
  $L_{\mathrm{eff}}/D$ is on a log scale, 0.01--1. Format PNG, \texttt{inferno}.
\item WHAT: 193M medians of $k/U_\infty^2$ are $1.4\times10^{-3}$, $1.3\times10^{-3}$, $2.4\times10^{-3}$, and of
  $L_{\mathrm{eff}}/D$ 0.07, 0.10, 0.14 over the strained points (0.15, 0.14, 0.16 over all points) at $x/D$ 4, 10,
  20. In the freestream $|S^*|\to0$, so $L_{\mathrm{eff}}=k^{1/2}|b|/|S^*|$ of the 3DNS is patchy, 0.05--1 (j--l);
  the RANS length there is uniform, set by its freestream $k$ and $\omega$.
\end{itemize}
\planfig{fig_cross_planes.png}{0.95}{F16 (\texttt{fig\_cross\_planes.py}): $k/U_\infty^2$ (a--f) and $L_{\mathrm{eff}}/D$
(g--l), RANS 50M and 193M 3DNS (b010), $x/D$ 4, 10, 20; white: $\eta$ 0.1--0.4. No masking: every point is drawn.
$L_{\mathrm{eff}}/D$ on a log scale, 0.01--1.}

\textbf{Figure F17: where length and alignment differ on the centre plane.} $\ln(L_{\mathrm{eff}}/L_{\mathrm{eff}}^{193})$
of the RANS 50M (top) and $\cos\theta$ of the 193M 3DNS (bottom), $x/D$ $-2$ to 20.
\begin{itemize}
\item WHY: locates the two parts of the F14 $L_{\mathrm{al}}$ term in the flow.
\item HOW: centre plane, (a) $\ln(L_{\mathrm{RANS}}/L_{193})$, (b) $\cos\theta$ of the 193M 3DNS (b010). No
  masking; only the wall row ($k=0$) is blank. Format PNG, \texttt{bwr}, symmetric limits.
\item WHAT: (a) the RANS length is short over the whole film; (b) the centre plane is the best-aligned part of the
  film (median $\cos\theta$ 0.83, against 0.62 and 0.72 over the film volume); misalignment sits in the jet, near
  the wall, in the hole and on the film flanks, which is where the alignment term of F14 comes from. Medians,
  points above the $|S^*|$ threshold / all points: $\ln$ ratio $+0.04/+0.22$ (jet), $-0.64/-0.35$ (near film),
  $-0.56/-0.42$ (far film); $\cos\theta$ 0.73/0.31, 0.83/0.71, 0.83/0.80; $\cos\theta<0$ at 18\,\% of all plane
  points (hole and freestream; none in the film). Film volume medians of $\cos\theta$: 0.62 (near), 0.72 (far).
\end{itemize}
\planfig{fig_centre_factors.png}{0.9}{F17 (\texttt{fig\_centre\_factors.py}): centre plane, (a)
$\ln(L_{\mathrm{RANS}}/L_{193})$, (b) $\cos\theta$ of the 193M 3DNS (b010). No masking; blank only the wall row
($k=0$).}

\textbf{Figure F18: how the turbulence budget and lengths vary along wall-normal lines.} Centreline and mid-pitch
at $x/D$ 4, 10, 20: (a) $P_k$ and $\varepsilon$ in units of $U_\infty^3/D$; (b) $L_{\mathrm{eff}}$ against
$C_\mu k^{3/2}/\varepsilon$; RANS 50M, 20M LES, 193M 3DNS.
\begin{itemize}
\item WHY: ties the energy and length factors to quantities the reader knows: where production sustains $k$, and
  that $L_{\mathrm{eff}}$ is the physical dissipation length in the film.
\item HOW: (a--c, d--f) $P_k$ (lines) and $\varepsilon$ (markers) over $\rho_\infty U_\infty^3/D$, centreline and
  mid-pitch at $x/D$ 4, 10, 20; (g--i, j--l) $L_{\mathrm{eff}}$ (lines) and $C_\mu k^{3/2}/\varepsilon$ (markers);
  193M (b010), 50M and 20M LES (b006), RANS 50M. The $k$ budget is E14.
\item WHAT: 193M medians over $z/D$ 0.05--1.5: $P_k/\varepsilon$ 1.27, 1.03, 0.90 (centreline) and 0.64, 1.57,
  1.34 (mid-pitch); $L_{\mathrm{eff}}/(C_\mu k^{3/2}/\varepsilon)$ 1.02--1.07 (centreline), 0.74--1.22
  (mid-pitch).
\end{itemize}
\planfig{fig_budget_lines.pdf}{0.95}{F18 (\texttt{fig\_budget\_lines.py}): (a--c, d--f) $P_k$ (lines) and
$\varepsilon$ (markers) over $\rho_\infty U_\infty^3/D$, centreline and mid-pitch at $x/D$ 4, 10, 20; (g--i, j--l)
$L_{\mathrm{eff}}$ (lines) and $C_\mu k^{3/2}/\varepsilon$ (markers); 193M (b010), 50M and 20M LES (b006), RANS
50M.}

\textbf{Figure F19: how strongly the reference diffuses heat sideways.} (a) Lateral over wall-normal eddy
diffusivity, $\nu_y/\nu_z$ and $\alpha_y/\alpha_z$, per region of the 193M b010; (b) lateral heat flux over the
film edges, RANS 50M against the 193M.
\begin{itemize}
\item WHY: step 3, one message: the reference heat diffusivity is anisotropic on the film flanks, as measured by
  Kaszeta and Simon (2000), and the RANS lateral diffusivity is about half, mostly through the $\Pr_t$ setting.
\item HOW: (a) $\nu_y/\nu_z$ and $\alpha_y/\alpha_z$ of the 193M b010 on the $x/D$ 10 and 20 cross-planes (all
  $z/D\ge0.3$, centre, flanks, wall layer), with the Kaszeta and Simon band 1.5--4 and the RANS value 1. (b) The
  lateral heat flux over the film edges at $x/D$ 10 and 20: $\ln(P_{\mathrm{RANS}}/P_{193})$ split exactly into
  $\mu_t$, anisotropy, $\Pr_t$ setting and gradient terms. The $\mu_t$ term compares the RANS $\mu_t$ with the
  classic wall-normal viscosity of the 193M, not the $\mu_{\mathrm{eff}}$ of F14.
\item WHAT: on the flanks $\alpha_y/\alpha_z$ is 3.2 and 1.7 and $\nu_y/\nu_z$ is 1.1 and 0.7. The split (b010) is
  $x/D=10$: $\mu_t$ $+0.91$, anisotropy $-0.80$, $\Pr_t$ $-0.72$, gradient $+0.66$, total $+0.04$; $x/D=20$:
  $+0.12$, $-0.11$, $-0.47$, $+0.65$, total $+0.19$.
\end{itemize}
\planfig{fig_prt_data.pdf}{0.8}{F19 (\texttt{fig\_prt\_data.py}): (a) $\nu_y/\nu_z$ and $\alpha_y/\alpha_z$ of the 193M
b010 on the $x/D$ 10 and 20 cross-planes (all $z/D\ge0.3$, centre, flanks, wall layer), Kaszeta and Simon band
1.5--4, RANS 1; (b) lateral heat-flux split over the film edges, RANS 50M against 193M.}

\textbf{(iii) What the literature offers against each factor} (end of 4.2, leads to 4.3).
\begin{itemize}
\item Length scale: no film-cooling-specific correction of the SST length; data-driven eddy viscosity (Ling et
  al., 2016b; Milani et al., 2020).
\item $\Pr_t$ too high: lower or variable $\Pr_t$ (Liu et al., 2008, 2011; Lakehal, 2002); data-driven diffusivity
  (Milani et al., 2018, 2020; Ellis and Xia, 2024; Laroche et al., 2023; Qaragoez and Mazaheri, 2025).
\item Spanwise and alignment: anisotropic eddy viscosity (Azzi and Lakehal, 2002); quadratic constitutive relation
  (Spalart, 2000; Sabnis et al., 2021).
\item Energy in the far film: the transition model holds the mid-pitch laminar; no film-cooling trip exists in
  $\gamma$--$\mathrm{Re}_\theta$ (Langtry and Menter, 2009).
\item Each remedy targets one factor, while the split shows all three in error, the length the most. Closing
  sentence: Section 4.3 tests which remedy recovers $\bar\eta$ a posteriori.
\end{itemize}

%-----------------------------------------------------------------------------------------------------
\subsection*{4.3\quad A Posteriori Tests}
\textbf{Purpose.} Section 4.2 names the closure inputs that differ from the reference. Section 4.3 changes one
input at a time in the RANS 50M and measures how much of the effectiveness error each change removes, and
whether the film it produces has the shape of the reference film. About 1.5 pages, one table and one figure.

\textbf{(i) Method.}
\begin{itemize}
\item WHY: a difference in a closure input matters only if correcting it changes $\eta$. The a priori terms of
  4.2 do not say how much each is worth.
\item HOW: every run starts from the converged RANS 50M and changes one input (Table T2). Four measures are
  reported against the 193M: the recovered share of the $\bar\eta$ deficit over $x/D$ 7--20; centreline $\eta$ at
  $x/D=20$; mid-pitch $\eta$ at $x/D=20$; $\bar\eta$ over the last part of the plate, $x/D$ 15--20. The first measures the amount
  of coolant at the wall; the other three measure where it is. BR is reported for each run (1.56--1.59 against
  1.527 for the 193M).
\item A run is counted as matching the reference only if the window mean and the three shape measures agree; a
  window mean alone can be met by a film of the wrong shape.
\end{itemize}

\textbf{Table T2: a posteriori runs.} One row per run: the input changed, where, the four measures and BR.
Purpose: the whole of 4.3 in one view. Results available:

{\centering\footnotesize
\begin{tabular}{@{}l l r r r r@{}}
\toprule
run & input changed & recovered & centreline $\eta$ & mid-pitch $\eta$ & $\bar\eta$ \\
 & & $x/D$ 7--20 & $x/D=20$ & $x/D=20$ & $x/D$ 15--20 \\
\midrule
SST, $\Pr_t$ 0.9 (base) & none & 0\,\% & 0.497 & 0.006 & 0.138 \\
$\Pr_t$ 0.5 / 0.4 / 0.3 & $\Pr_t$, whole domain & 36 / 44 / 48\,\% & 0.460 / 0.422 / 0.361 & 0.009 / 0.010 / 0.012 & 0.160 / 0.164 / 0.164 \\
$\Pr_t$ 0.4 near the hole & $\Pr_t$, $x/D<7.5$ & 34\,\% & 0.469 & 0.008 & 0.157 \\
QCR & stress relation & 72\,\% & 0.524 & 0.006 & 0.185 \\
QCR, $\Pr_t$ 0.4 & both & 97\,\% & 0.426 & 0.010 & 0.198 \\
QCR, $\Pr_t$ 0.4 near the hole & both & 96\,\% & 0.475 & 0.008 & 0.200 \\
3DNS heat flux & heat flux, all regions & 76\,\% & 0.339 & 0.036 & 0.191 \\
$\Pr_t$ 0.4 in the film only & $\Pr_t$, $x/D\ge3$ & & & & \\
3DNS heat flux in the film only & heat flux, $x/D\ge3$ & & 0.360 & 0.031 & 0.183 \\
3DNS stresses & stress, all regions & & & & \\
3DNS stresses and heat flux & both, all regions & & 0.227 & 0.035 & 0.178 \\
$\gamma=1$ downstream of the hole & transition model & & & & \\
\midrule
193M 3DNS & & & 0.387 & 0.025 & 0.208 \\
experiment & & & 0.419 & 0.035 & 0.212 \\
\bottomrule
\end{tabular}\par}
Notes for the table: the imposed stresses and heat fluxes are those of the 50M 3DNS; the run with both imposed
has BR 1.416 and a different approach layer, so its centreline deficit is not a closure result alone.

\textbf{(ii) The $\Pr_t$ setting} (step 2 of 4.2).
\begin{itemize}
\item WHY: the reference $\Pr_t$ is 0.23--0.47 against the setting of 0.9. The test is whether lowering the
  setting recovers $\eta$.
\item HOW: uniform $\Pr_t$ of 0.5, 0.4 and 0.3; $\Pr_t$ 0.4 near the hole only; $\Pr_t$ 0.4 in the film only.
\item WHAT: the recovery saturates at 36--48\,\%. The centreline $\eta$ passes through the reference value
  between $\Pr_t$ 0.4 and 0.3, while $\bar\eta$ over $x/D$ 15--20 stays at 0.160--0.164 (reference 0.208) and the mid-pitch $\eta$ at
  0.009--0.012. A lower setting moves coolant off the centreline and does not carry it to mid-pitch.
\item Film-only run, likely outcomes: (a) recovery of 35--45\,\%, as for the uniform 0.4: the film setting carries
  the effect and the jet setting is immaterial; (b) below 30\,\%: the effect is set in the jet, where the
  reference $\Pr_t$ is lowest; (c) above 45\,\%: lowering $\Pr_t$ in the jet is harmful. The near-hole run (34\,\%
  from $x/D<7.5$ alone) makes (a) or (b) the more likely.
\end{itemize}

\textbf{(iii) The form of the heat flux} (steps 2 and 3).
\begin{itemize}
\item WHY: the reference heat diffusivity is anisotropic on the film flanks, which no $\Pr_t$ setting gives.
\item HOW: the heat flux of the 50M 3DNS is imposed in the RANS energy equation and the SST stresses are kept;
  once in all regions and once in the film only. The imposed stresses of (iv) are also those of the 50M 3DNS.
\item WHAT: recovery 76\,\%; the only run with coolant at mid-pitch (0.036 against 0.025); $\bar\eta$ over $x/D$ 15--20
  0.191 against 0.208. The centreline is below the reference (0.339 against 0.387): with the RANS mean flow
  unchanged, the imposed flux spreads the coolant further than in the reference. The window mean overstates the
  match.
\item Film-only run: centreline 0.360 against 0.387, mid-pitch 0.031 against 0.025, $\bar\eta$ over $x/D$ 15--20
  0.183 against 0.208, and a lateral profile that follows the reference to $|y/D|=2.9$. It is the closest run
  on all shape measures; the over-spreading of the all-region run is reduced, so part of that came from the jet.
  The mean flow of the RANS is unchanged in both runs: correcting the turbulent heat flux in the film recovers
  most of the film shape with the too-strong CRVP of 4.1 still in place.
\end{itemize}

\textbf{(iv) The form of the stress} (the part outside the split of 4.2).
\begin{itemize}
\item WHY: the eddy viscosity of 4.2 carries only the part of the stress aligned with the strain. The quadratic
  constitutive relation adds a part with zero contraction with the strain, so it leaves $\mu_{\mathrm{eff}}$
  unchanged; it tests what the split cannot see.
\item HOW: QCR in the whole domain; the stresses of the 3DNS imposed in the momentum equations; the QCR stress
  decomposed into its aligned and perpendicular parts, and the change of the mean secondary flow measured.
\item WHAT: QCR recovers 72\,\%, raises $\bar\eta$ over $x/D$ 15--20 from 0.138 to 0.185 and raises the centreline
  (0.524 against 0.497 for the base and 0.387 for the reference); the mid-pitch $\eta$ is unchanged. QCR with
  $\Pr_t$ 0.4 recovers 97\,\% and brings the centreline to 0.426. The two shares add to more than the combined
  run (72 + 44 against 97), so the corrections overlap.
\item QCR decomposition, likely outcomes: (a) regional $\mu_t$ within about 10\,\% of the base and a changed
  streamwise vorticity: QCR acts through the secondary flow driven by the normal-stress differences; (b) regional
  $\mu_t$ changed by more: QCR acts in part through the eddy viscosity, by changing the mean strain; (c) both.
\item 3DNS-stress run, likely outcomes: (a) recovery of 60--90\,\%, at or above QCR: the stress, aligned and
  perpendicular parts together, supplies most of the deficit; (b) below QCR: part of the QCR gain is a
  cancellation and is not supported by the reference stress; (c) no converged solution: reported as such.
\item Combined run, likely outcomes: (a) 85--110\,\% with all shape measures within the reference: the two closures
  account for the error; (b) high recovery with the mid-pitch still low: the laminar mid-pitch layer is a
  separate error; (c) the parts are sub-additive, stress plus heat flux exceeding the combined run, as for QCR and
  $\Pr_t$.
\end{itemize}

\textbf{(v) The laminar mid-pitch layer} (the energy factor).
\begin{itemize}
\item WHY: the RANS layer between the jets stays laminar to the end of the plate, where the reference is
  intermittent and then turbulent (4.1).
\item HOW: the intermittency is set to 1 downstream of the hole exit; approach layer and hole unchanged.
\item WHAT, likely outcomes: (a) recovery below 20\,\% and mid-pitch $\eta$ still well below 0.025: a turbulent
  layer alone does not bring coolant to mid-pitch, and the lateral transport is the limit (the plate-mean $\eta$
  rises by 1.3\,\% on the 5M mesh, which points here); (b) recovery of 20--60\,\%: the missing turbulence between
  the jets is a comparable error to the $\Pr_t$ setting; (c) $\bar\eta$ falls: more mixing thins the film.
\end{itemize}

\textbf{(vi) What is not tested.}
\begin{itemize}
\item No run changes the turbulence length scale alone. The length of 4.2 is tested only as part of the 3DNS-stress
  run, which replaces the whole eddy viscosity.
\item BR differs between the runs by up to 2\,\% (QCR 1.561, base 1.593); from the measured BR sensitivity (4.1)
  this moves $\bar\eta$ by much less than the recoveries reported.
\end{itemize}

\textbf{Figure F20: the film of each corrected run against the reference.} Four panels, key runs (base,
$\Pr_t$ 0.4, QCR with $\Pr_t$ 0.4, 3DNS heat flux in the film only; $\gamma=1$ and the zonal $\Pr_t$ run when
complete) against the 193M and the experiment. The other runs are in Table T2 only.
\begin{itemize}
\item WHY: the table gives single numbers; the figure shows whether a run recovers the film or only its mean.
\item HOW: (a) $\bar\eta$, (b) centreline $\eta$, (c) mid-pitch $\eta$ against $x/D$; (d) the lateral
  profile $\eta(|y|)$ averaged over $x/D$ 15--20. 193M as a thick line, the experiment (DR 1.5, BR 1.5, Tu
  0.5\,\%) as markers, each RANS run with its own colour and line style.
\item WHAT: a lower $\Pr_t$, with or without QCR, brings the centreline to within 0.04 of the reference and leaves
  the mid-pitch at 0.010 and the lateral profile narrow; the imposed film heat flux follows the reference in the
  centreline and the lateral profile and over-shoots the mid-pitch from $x/D=8$. The measured mid-pitch $\eta$ is
  flat at about 0.035 from $x/D=2$, above every simulation: the laminar approach layer of the configuration
  (4.1 (iv)), referred to in the caption.
\end{itemize}
\planfig{fig_aposteriori.pdf}{0.95}{F20 (\texttt{fig\_aposteriori.py}): (a) $\bar\eta$, (b) centreline and (c)
mid-pitch $\eta$ against $x/D$, (d) lateral profile of $\eta$ over $x/D$ 15--20; corrected RANS 50M runs, the
193M 3DNS and the experiment.}

%=====================================================================================================
\section*{5\quad Conclusions}
Section 5 will contain the conclusions, which mirror the aims of Section 2.

\end{document}
